\DocumentMetadata{lang=en-US,tagging=on}
\documentclass[11pt,a4paper]{article}
\usepackage{fontspec}
\setmainfont{TeX Gyre Pagella}
\setsansfont{TeX Gyre Heros}
\usepackage{microtype}
\usepackage[margin=27mm,headheight=15pt]{geometry}
\usepackage{amsmath}
\usepackage{unicode-math}
\setmathfont{TeX Gyre Pagella Math}
\newcommand{\square}{\mdlgwhtsquare}
\usepackage{aliascnt}
\usepackage{booktabs,longtable,array,tabularx}
\usepackage{xcolor,graphicx,tikz}
\usetikzlibrary{arrows.meta,positioning,fit,calc}
\usepackage[most]{tcolorbox}

\usepackage{needspace}
\usepackage{float,algorithm,algpseudocode}
\usepackage{fancyhdr}
\usepackage{etoolbox}
\usepackage{hyperref}
\usepackage[nameinlink,noabbrev]{cleveref}
\definecolor{ink}{HTML}{173449}
\definecolor{accent}{HTML}{176D83}
\definecolor{pale}{HTML}{EDF5F7}
\definecolor{muted}{HTML}{546574}
\definecolor{typeorange}{HTML}{A34B00}
\definecolor{typeblue}{HTML}{185F96}
\newcommand{\otype}{\textcolor{typeorange}{\mathsf O}}
\newcommand{\btype}{\textcolor{typeblue}{\mathsf B}}
\hypersetup{colorlinks=true,linkcolor=accent,citecolor=accent,urlcolor=accent,
 pdftitle={Information, Belief, and Search},pdfauthor={MTGallium},
 pdfsubject={A formal reference and guided introduction to imperfect-information game search}}
\urlstyle{same}
\makeatletter
\renewcommand\section{\@startsection{section}{1}{\z@}
  {-3.5ex \@plus -1ex \@minus -.2ex}{2.3ex \@plus .2ex}
  {\normalfont\Large\sffamily\bfseries\color{ink}}}
\renewcommand\subsection{\@startsection{subsection}{2}{\z@}
  {-3.25ex \@plus -1ex \@minus -.2ex}{1.5ex \@plus .2ex}
  {\normalfont\large\sffamily\bfseries\color{ink}}}
\makeatother
\tagpdfsetup{table/header-rows={1}}
\pretocmd{\section}{\Needspace{11\baselineskip}}{}{}
\pretocmd{\subsection}{\Needspace{9\baselineskip}}{}{}
\setlength{\parindent}{0pt}
\setlength{\parskip}{5pt plus 1pt}

\setlength{\emergencystretch}{2em}
\pagestyle{fancy}
\fancyhf{}
\fancyhead[L]{\small\sffamily\color{muted}MTGallium / Formal foundations}
\fancyhead[R]{\small\sffamily\color{muted}Information, Belief, and Search}
\fancyfoot[C]{\small\thepage}
\renewcommand{\headrulewidth}{.3pt}
\newcounter{definition}[section]
\renewcommand{\thedefinition}{\thesection.\arabic{definition}}
\crefname{definition}{definition}{definitions}
\NewDocumentEnvironment{definition}{o}{\par\refstepcounter{definition}\textbf{Definition \thedefinition}\IfValueT{#1}{ (#1)}. \itshape}{\par}
\newaliascnt{proposition}{definition}
\crefname{proposition}{proposition}{propositions}
\NewDocumentEnvironment{proposition}{o}{\par\refstepcounter{proposition}\textbf{Proposition \theproposition}\IfValueT{#1}{ (#1)}. \itshape}{\par}
\newaliascnt{lemma}{definition}
\crefname{lemma}{lemma}{lemmas}
\NewDocumentEnvironment{lemma}{o}{\par\refstepcounter{lemma}\textbf{Lemma \thelemma}\IfValueT{#1}{ (#1)}. \itshape}{\par}
\newaliascnt{example}{definition}
\crefname{example}{example}{examples}
\NewDocumentEnvironment{example}{o}{\par\refstepcounter{example}\textbf{Example \theexample}\IfValueT{#1}{ (#1)}. }{\par}
\newaliascnt{remark}{definition}
\crefname{remark}{remark}{remarks}
\NewDocumentEnvironment{remark}{o}{\par\refstepcounter{remark}\textbf{Remark \theremark}\IfValueT{#1}{ (#1)}. }{\par}
\NewDocumentEnvironment{proof}{O{Proof}}{\par\textit{#1.} }{\hfill\fbox{\rule{0pt}{3pt}\rule{3pt}{0pt}}\par}
\newtcolorbox{principle}[1]{breakable,colback=pale,colframe=accent,
 boxrule=.5pt,arc=1mm,title={#1},fonttitle=\sffamily\bfseries,
 colbacktitle=accent,coltitle=white,left=3mm,right=3mm,top=2mm,bottom=2mm}
\newcommand{\Hist}{\mathcal H}
\newcommand{\Info}{\mathcal I}
\newcommand{\Obs}{\Omega}
\newcommand{\Act}{\mathcal A}
\newcommand{\Init}{S_{\mathrm{init}}}
\newcommand{\last}{\operatorname{last}}
\newcommand{\nexts}{\operatorname{next}}
\newcommand{\own}{\operatorname{own}}
\newcommand{\cat}{\mathbin{\Vert}}
\newcommand{\epss}{\mathsf E}
\newcommand{\Prob}{\mathbb P}
\newcommand{\Expect}{\mathbb E}
\newcommand{\ind}{\mathbf 1}
\newcommand{\supp}{\operatorname{supp}}
\newcommand{\push}{_{\#}}
\newcommand{\code}[1]{\texttt{\small #1}}
\newcommand{\contract}[1]{\textbf{Contract.} #1}
\newcommand{\sourcepath}[2]{\href{https://github.com/huxinq/MTGallium/blob/main/#1}{#2}}

\begin{document}
\hypersetup{pageanchor=false}
\begin{titlepage}
\vspace*{12mm}
{\sffamily\color{accent}\large MTGALLIUM WHITE PAPER\par}
\vspace{12mm}
{\sffamily\bfseries\color{ink}\fontsize{34}{39}\selectfont
Information, Belief,\\and Search\par}
\vspace{7mm}
{\Large A formal reference and guided introduction\\to imperfect-information game search\par}
\vspace{12mm}
\begin{tikzpicture}[alt={Four-part reading map. Game and information leads to probability and targets, then simulation and search, then representations and implementation.},node distance=8mm and 9mm,
 every node/.style={font=\small\sffamily},
 block/.style={draw=accent,fill=pale,rounded corners=2pt,
 text width=35mm,minimum height=13mm,align=center},
 arr/.style={-{Latex[length=2mm]},thick,draw=accent}]
\node[block,text width=58mm] (g) {I. Game and information\\What happens; what each player observes};
\node[block,text width=58mm,right=of g] (h) {II. Probability and targets\\What is plausible; what a choice is worth};
\node[block,text width=58mm,below=of h] (s) {III. Simulation and search\\How hypothetical play produces estimates};
\node[block,text width=58mm,left=of s] (r) {IV. Representations and implementation\\What the computer must preserve};
\draw[arr] (g)--(h); \draw[arr] (h)--(s); \draw[arr] (s)--(r);
\end{tikzpicture}
\vfill
\begin{minipage}{.92\textwidth}
\textbf{Purpose.} A reference and a teaching document: definitions, invariants,
short proofs, worked examples, and the conditions that connect theoretical
objects to a practical game-search implementation.

\textbf{Scope.} Two players; deterministic engine advancement; potentially
randomized player choices; imperfect information; full observation histories.
\end{minipage}\par
\vspace{10mm}
\noindent\begin{minipage}{\textwidth}
\sffamily\textbf{Version 1.4}\quad / \quad 9 September 2026\\[2mm]
\small A guided development with worked examples and a formal reference.
\end{minipage}\par
\end{titlepage}
\pagenumbering{roman}
\hypersetup{pageanchor=true}

\begin{abstract}
A player sees a sequence of observations and must choose an action without
knowing the complete game history. A signal game lets us follow that decision:
observations narrow the possibilities, a model of the sender weights them,
and a rule for future choices determines expected payoffs. Worked updates and
a search trace show how those probabilities become action estimates. We also
show why a search can use only the root player's information yet give its
simulated opponent an information advantage: two sender-root searches can
produce different guesses from the same opponent observations. The final part
asks what an implementation must retain to reproduce these decisions and
values, and connects the objects to MTGallium source. Definitions, proofs,
exercises, and separate solutions provide a reference alongside the examples.
\end{abstract}

\begingroup
\small
\setlength{\parskip}{0pt}
\makeatletter
\let\originaltocsection\l@section
\renewcommand{\l@section}[2]{\vspace{-3pt}\originaltocsection{#1}{#2}}
\makeatother
\tableofcontents
\endgroup
\clearpage
\pagenumbering{arabic}

\section{Begin with a private signal}
\label{sec:reading}

A sender knows a private type, \textcolor{typeorange}{orange} or \textcolor{typeblue}{blue}. The sender chooses a public signal,
and a second player tries to guess the type. The same signal can come from
either type. What should the guesser do, and how could a search program decide?

The game is small, but answering those questions calls for the same objects
that a card-game planner uses. A game history tells us which type the sender
actually has and which signal was chosen. The guesser observes only part of
that history. A belief assigns probabilities to the remaining possibilities.
Those probabilities, together with a rule for future choices, determine expected
payoffs. Search estimates the payoffs by trying hypothetical continuations.

We will return to this game as each object becomes available. Sender is player
$0$, guesser is player $1$, and the hidden type is $\theta\in\{\otype,\btype\}$. The
signals are $\ell$ (left) and $\mathsf r$ (right); the guesses are $g_{\otype}$ and $g_{\btype}$. A correct guess wins for the guesser;
an incorrect guess wins for the sender. The type is fixed before play. Only
player choices may randomize.

\begin{principle}{The question that guides the construction}
After seeing $\ell$, the guesser must choose the same distribution over guesses
in the orange and blue cases whenever its observations are identical. Search can
examine both cases to estimate their consequences. The challenge is to use
those hypotheses to choose well with the information the player actually has.
\end{principle}

\subsection{A reading map}

\textbf{Part I} builds game histories and projects them into each player's
experience. It then collects the histories compatible with that experience.
\textbf{Part II} adds policies and probabilities, leading to a concrete
calculation of what the two guesses are worth. \textbf{Part III} follows search
iterations and their statistics, including the different ways to model a
simulated opponent. \textbf{Part IV} gives conditions under which the entire
computation preserves the root player's information boundary and connects the
objects to MTGallium's representations.

\Cref{app:notation,app:contracts} collect notation and classified contracts;
\cref{app:exercises,app:solutions} provide exercises and solutions.

This paper brings together standard ideas from extensive-form games, partially
observable models, and Monte Carlo search to specify the information and value
semantics of MTGallium. \Cref{sec:related} gives the correspondence to the
literature; \cref{sec:implementation} maps the objects to source.

\subsection{Three kinds of statement}

\emph{Model assumptions} fix the game we study: for example, automatic processing
terminates. \emph{Mathematical consequences} follow from those assumptions and
definitions: for example, indistinguishable decision histories have the same
legal menu. \emph{Implementation obligations and algorithm choices} explain
how a program realizes the objects and allocates computation. This classification
will also organize the reference sheet.

We use finite action sets and finite or countable state and observation spaces,
so probabilities can be written as sums.

\clearpage\part{Game and information}

\section{Game states, choices, and advancement}
\label{sec:game}

A \emph{game state} is a state awaiting a player choice or a terminal state.
Advancement collects the automatic rule processing between these points into
one transition. An engine can retain finer intermediate configurations for
execution and diagnosis.

\Needspace{7\baselineskip}
\begin{definition}[Deterministic game skeleton]
A game skeleton consists of
\[
\mathcal G_{\mathrm{dyn}}=(P,S,\Init,F,\Act,p,A,\nexts),
\]
where $P=\{0,1\}$ is the player set, $S$ is the game-state space,
$\varnothing\ne\Init\subseteq S$ is the allowed initial-state set,
$F\subseteq S$ is the terminal-state set, and $\Act$ is an action universe.
For every nonterminal state $s$,
\[
p(s)\in P,\qquad \varnothing\ne A(s)\subseteq\Act,
\]
with $A(s)$ finite. Advancement is a function
\[
\nexts:\{(s,a):s\notin F,\ a\in A(s)\}\longrightarrow S.
\]
\end{definition}

The action $a$ includes the semantic content of the choice. If a target or an
ordering requires another player choice, advancement stops at that decision. If a game presents an entire compound choice at once,
the action can encode that complete choice. Granularity follows the actual
decision interface rather than an arbitrary size of engine instruction.

\contract{For each legal $(s,a)$, advancement consumes that action and all
automatic consequences, then reaches the next player choice or termination
after finitely many automatic steps.}

\subsection{Determinism and hidden information}

Given the same $s$ and $a$, the same next state results. The acting player need
not alternate: $p(\nexts(s,a))$ can equal $p(s)$. Terminal states have neither
an acting player nor an outgoing action transition. We leave $p,A$ undefined
there.

Drawing the next card from a fixed hidden library order is deterministic: the
order in $s$ determines the card, even when the player is uncertain about it.
A belief or prior will describe that uncertainty. Fresh random shuffles require
a chance-transition extension, outside the present model. Here randomness enters
through player policies and through the distribution used to describe possible
initial states.

\begin{example}[Automatic consequence versus additional choice]
A player chooses to cast a spell. Paying a fully specified cost and putting the
spell on the stack may be automatic consequences of that action. Asking the
other player whether to respond is a new game state. Resolving an effect that
asks someone to choose a card likewise reaches a new game state. These
decision points determine the granularity of the history.
\end{example}

\subsection{Histories and reconstruction}

\Needspace{7\baselineskip}
\begin{definition}[Game history]
A valid finite game history is
\[
h_n=(s_0,a_0,s_1,\ldots,a_{n-1},s_n),
\]
where $s_0\in\Init$, and for $0\le j<n$, $s_j\notin F$,
$a_j\in A(s_j)$, and $s_{j+1}=\nexts(s_j,a_j)$.
Write $\Hist$ for all such histories, $|h_n|=n$, and $\last(h_n)=s_n$.
\end{definition}

Write $h\preceq h'$ when $h$ is a prefix of $h'$. Given $s_0$ and a legal
action sequence, the game history is uniquely reconstructible by repeated
application of $\nexts$. This reconstructs all game states at our chosen
granularity. It does not select a unique low-level engine trace unless an
implementation separately specifies that trace.

\begin{proposition}[History reconstruction]
Two valid histories with the same initial state and the same action sequence
have the same state sequence.
\end{proposition}
\begin{proof}
Their first states agree. If their states at index $j$ agree, then the common
action $a_j$ and the function $\nexts$ give the same state at index $j+1$.
Induction completes the argument.
\end{proof}

An unbounded sequence of individually terminating choices remains possible.
\Cref{sec:value} supplies the additional termination assumption needed for
terminal expected values.

\subsection*{The signal game}

Initial states $s^\theta$ await the sender's choice:
\[
p(s^\theta)=0,\qquad A(s^\theta)=\{\ell,\mathsf r\}.
\]
After a signal $m$, state $t^{\theta,m}$ awaits a guess:
\[
\nexts(s^\theta,m)=t^{\theta,m},\qquad
p(t^{\theta,m})=1,\qquad A(t^{\theta,m})=\{g_{\otype},g_{\btype}\}.
\]

A guess advances to terminal state $f^{\theta,m,g}$. Thus each complete play has two
transitions, one for the signal and one for the guess.

\section{Observations and player history}
\label{sec:observation}

Observation rules specify who learns what during play. A private reveal
can disclose a card identity without changing the eventual game configuration.
Conversely, an opponent can change a hidden configuration without disclosing
its new contents. Observation rules are additional components of the game.

\Needspace{7\baselineskip}
\begin{definition}[Observation functions]
For each player $q\in P$, let $\Obs_q$ be an observation alphabet and
$\Obs_q^*$ its finite words, including the empty word $\varepsilon$.
Specify
\[
O_q^0:\Init\to\Obs_q^*,\qquad
O_q:\operatorname{dom}(\nexts)\to\Obs_q^*.
\]
$O_q^0(s_0)$ is the initial observation sequence. $O_q(s,a)$ is the ordered
sequence disclosed during advancement by $a$, including any resulting request
for a choice. Concatenation is written $\cat$.
\end{definition}

One transition may contribute several observations. An empty contribution
leaves the observation word unchanged, preserving the privacy of an undisclosed
transition. Timing, turn, and phase information enter this word precisely when
the game makes them available to the player.

\contract{The chosen game state contains all context needed to determine both
advancement and its disclosures. In particular, equal $(s,a)$ pairs produce
equal $O_q(s,a)$ for each player.}

This assumption makes the observation function usable from a current
configuration. For example, a permission to inspect a card belongs in that
configuration when it affects a forthcoming disclosure. A history-dependent
disclosure rule requires either retaining the relevant context in the state
or giving the observation function a history argument. We use the former
formulation throughout.

\Needspace{7\baselineskip}
\begin{definition}[Player history and information state]
For $h_n\in\Hist$, define
\begin{equation}
I_q(h_n)=O_q^0(s_0)\cat O_q(s_0,a_0)\cat\cdots\cat
O_q(s_{n-1},a_{n-1}).
\label{eq:information}
\end{equation}
The \emph{player history} and the theoretical \emph{information state} are this
same object. Let $\Info_q=\{I_q(h):h\in\Hist\}$ be the realizable information
states. A current snapshot can summarize part of this accumulated experience.
\end{definition}

\begin{example}[Disclosed information and remembered consequence]
Earlier, a player observed ``the top library card is Lightning Bolt.'' Later,
they observe ``the opponent draws that top card.'' The second observation need
not repeat the card identity. Together the observations entail a remembered
fact about the opponent's hand, provided no intervening event invalidated it.
The observation function supplies disclosures; reasoning over the history
supplies their consequences.
\end{example}

\subsection*{The signal game}

For readability, write $\operatorname{request}(M)$ for a request to choose from
menu $M$, $\operatorname{chose}(a)$ for one's own realized choice, and
$\operatorname{signal}(m)$ for an observed public signal. These are typed
observation tokens, not engine identifiers.

The sender initially receives its type and a signal request. The guesser
initially receives its role. On the signal transition, the guesser receives
the public signal and a request to guess. In particular,
\begin{align*}
I_0(s^\theta)&=(\operatorname{type}(\theta),
                  \operatorname{request}(\{\ell,\mathsf r\})),\\
I_1(s^\theta,\ell,t^{\theta,\ell})
 &= (\operatorname{role}(\text{guesser}),\operatorname{signal}(\ell),
     \operatorname{request}(\{g_{\otype},g_{\btype}\}))=:i_\ell.
\end{align*}
The second line is identical for $\theta=\otype$ and $\theta=\btype$. On a guess
transition the guesser receives its own-choice token and the outcome. The
sender likewise receives the outcome after its earlier own-signal token.
These records explain both who must act and what that player remembers.

\subsection{Recovering one's own choices}

Define the action subsequence
\[
\own_q(h_n)=(a_j:0\le j<n,\ p(s_j)=q).
\]
Our own-choice requirement is the existence of
\begin{equation}
R_q:\Info_q\to\Act^*,\qquad R_q(I_q(h))=\own_q(h).
\label{eq:own}
\end{equation}
Equivalently, equal player histories imply equal own-action subsequences.
This applies to every valid prefix, not just completed games.

Perfect recall includes what the player knew when choosing each action.
Place each own-choice token before the observations generated by that choice:
when $p(s)=q$, let $O_q(s,a)$ begin with the typed token
$\operatorname{chose}_q(a)$, followed by the resulting disclosures. Initial
and unrelated observations use other token types. The prefix before each
own-choice token then records exactly the information available when the
action was selected, so one can recover
\[
\bigl((I_q(h_j),a_j):p(s_j)=q\bigr),
\]
together with the order of observations between choices. Other encodings can
preserve the same chronology.

\subsection{Recovering decision status and legal choices}

For every state, define
\[
d_q(s)=
\begin{cases}
A(s),&s\notin F\text{ and }p(s)=q,\\
\bot,&\text{otherwise}.
\end{cases}
\]
The symbol $\bot$ means that no choice is currently required from $q$. It does
not distinguish an opponent decision from termination. If terminal status is
itself disclosed, that is supplied by the observation rules.

Require a decoder
\begin{equation}
D_q:\Info_q\to\mathcal P_{\mathrm{fin}}(\Act)\cup\{\bot\},\qquad
D_q(I_q(h))=d_q(\last(h)).
\label{eq:decision}
\end{equation}
Only nonempty finite sets occur as decision values. A concrete implementation
need not print the full menu, but the theoretical history must determine it.
Typed request and action tokens are one sufficient encoding.

\begin{proposition}[Decision consistency]
If $I_q(h)=I_q(h')$, then $d_q(\last(h))=d_q(\last(h'))$.
In particular, if one is a decision for $q$, both are, and their legal action
sets are equal.
\end{proposition}
\begin{proof}
Apply the same function $D_q$ to the common information state.
\end{proof}

When two situations have different legal menus, check whether the player can
observe the difference. An observable menu difference belongs in the information
state. If the player receives no distinguishing observation, either the action
abstraction or the observation model needs revision.

\section{Compatibility and epistemic state}
\label{sec:compatibility}

\Needspace{7\baselineskip}
\begin{definition}[Compatible histories]
For realizable $i_q\in\Info_q$, define
\[
\mathcal C_q(i_q)=\{h\in\Hist:I_q(h)=i_q\}.
\]
These are the possibilities allowed by the game and observation rules.
\end{definition}

Define $h\sim_q h'$ exactly when $I_q(h)=I_q(h')$. Equality makes this
relation reflexive, symmetric, and transitive. Its equivalence classes are the
sets $\mathcal C_q(i_q)$; two such sets are identical or disjoint.
The actual history is always included:
\[
h\in\mathcal C_q(I_q(h)).
\]
Own-choice and decision consistency hold throughout each class.

\subsection{Three different compatible sets}

Define the full epistemic state of a history by
\begin{equation}
\epss(h)=(\last(h),I_0(h),I_1(h)).
\label{eq:epistemic}
\end{equation}
``Epistemic'' means concerning information or knowledge. The full tuple records
the current configuration together with both players' experiences. Its
continuation role is established below. The compatible current-state sets are
\begin{align}
\mathcal S_q(i_q)&=\{\last(h):h\in\mathcal C_q(i_q)\},\\
\mathcal E_q(i_q)&=\{\epss(h):h\in\mathcal C_q(i_q)\}.
\end{align}
The player information-state component $i_q$ is fixed in the second set; the
other player's component need not be. Both sets are images of a compatibility
class: applying each projection gathers its possible current tuples or states.

\begin{figure}[htbp]
\centering
\begin{tikzpicture}[alt={Four histories share the same player information. The full epistemic tuple groups histories one and two together and keeps three and four distinct. Current game state merges one, two, and three; four remains distinct.},x=1cm,y=1cm,every node/.style={font=\small},
 group/.style={draw=accent,fill=pale,rounded corners,minimum height=9mm},
 arr/.style={-{Latex},draw=muted}]
\node[anchor=east] at (0,2.1) {Histories};
\node (h1) at (1,2.1) {$h_1$}; \node (h2) at (2.5,2.1) {$h_2$};
\node (h3) at (4.5,2.1) {$h_3$}; \node (h4) at (6.5,2.1) {$h_4$};
\node[anchor=east] at (0,1.3) {$I_q$ partition};
\node[group,minimum width=59mm] at (3.75,1.3) {$\{h_1,h_2,h_3,h_4\}$};
\node[anchor=east] at (0,.2) {$\mathsf E$ partition};
\node[group,minimum width=23mm] (e1) at (1.75,.2) {$\{h_1,h_2\}$};
\node[group,minimum width=16mm] (e2) at (4.5,.2) {$\{h_3\}$};
\node[group,minimum width=16mm] (e3) at (6.5,.2) {$\{h_4\}$};
\node[anchor=east] at (0,-1) {$\operatorname{last}$ partition};
\node[group,minimum width=48mm] (s1) at (2.8,-1) {$\{h_1,h_2,h_3\}$};
\node[group,minimum width=16mm] (s2) at (6.5,-1) {$\{h_4\}$};
\draw[arr] (e1)--(s1); \draw[arr] (e2)--(s1); \draw[arr] (e3)--(s2);
\end{tikzpicture}
\caption{An illustrative family with one fixed $I_q$. Histories $h_1,h_2$
differ only in discarded past detail and share the full tuple. History $h_3$
has the same current game state but different opponent memory; $h_4$ has a
different current state. Thus this family yields three epistemic tuples and
two game states. Globally, game-state and player-information partitions can
cross; the full tuple retains both.}
\end{figure}

\begin{example}[The same game state, different remembered information]
Two histories end with the same cards in the same zones, the same resources,
and the same player to act. In one, the opponent earlier observed a particular
card in the player's hand; in the other, they did not. If the earlier reveal
leaves no current rule-relevant flag, the final game states can agree while
the opponent's information histories differ. The full epistemic states remain
distinct even though their first components coincide.
\end{example}

Compatibility is joint. A game state from $\mathcal S_q(i_q)$ and an opponent
information state occurring somewhere in $\mathcal E_q(i_q)$ cannot necessarily
be combined. A valid tuple needs a common witnessing history. Independently
sampling these components can fabricate a situation that never occurred in
any compatible play.

\subsection*{The guesser after the left signal}

After the left signal, the guesser has information state $i_\ell$. Exactly two
histories in the basic game produce it:
\[
\mathcal C_1(i_\ell)=\{(s^{\otype},\ell,t^{\otype,\ell}),\ (s^{\btype},\ell,t^{\btype,\ell})\}.
\]
Their current game states differ in private type. Their full epistemic tuples
also differ because the sender remembers that type. The guesser's component
is the same. A single guess therefore has two possible consequences.

\begin{figure}[htbp]
\centering
\begin{tikzpicture}[alt={After the left signal, orange-type and blue-type histories share one guesser information-state node. Guessing orange wins in the orange hypothesis and loses in the blue hypothesis.},x=1cm,y=1cm,
 state/.style={draw=accent,rounded corners,fill=pale,minimum width=22mm,
 minimum height=9mm,font=\small},
 arr/.style={-{Latex},thick,draw=accent}]
\node[state,draw=typeorange,fill=typeorange!6] (r) at (0,1) {$s^{\otype}$};
\node[state,draw=typeblue,fill=typeblue!6] (b) at (0,-1) {$s^{\btype}$};
\node[state,draw=typeorange,fill=typeorange!6] (tr) at (4,1) {$t^{\otype,\ell}$};
\node[state,draw=typeblue,fill=typeblue!6] (tb) at (4,-1) {$t^{\btype,\ell}$};
\node[state,draw=typeorange,fill=typeorange!6] (fr) at (8,1) {$+1$ to guesser};
\node[state,draw=typeblue,fill=typeblue!6] (fb) at (8,-1) {$-1$ to guesser};
\draw[arr,draw=typeorange] (r)--node[above] {$\ell$}(tr);
\draw[arr,draw=typeblue] (b)--node[above] {$\ell$}(tb);
\draw[arr,draw=typeorange] (tr)--node[above] {$g_{\otype}$}(fr);
\draw[arr,draw=typeblue] (tb)--node[above] {$g_{\otype}$}(fb);
\node[draw=muted,dashed,rounded corners,fit=(tr)(tb),inner sep=5mm,
 label={[font=\small\sffamily]above:One information-state node $n_\ell$}] {};
\end{tikzpicture}
\caption{Two concrete trajectories share a guesser decision node. Choosing $g_{\otype}$
has different outcomes in its two compatible hypotheses. Only the $\ell$ branch is shown.}
\end{figure}

\subsection{Remembering facts without assigning probabilities}

For a predicate $\varphi$ on game histories, say it is entailed by $i_q$ when
\[
\forall h\in\mathcal C_q(i_q),\quad\varphi(h).
\]
For a current-game-state predicate $\psi$, use
$\forall s\in\mathcal S_q(i_q),\ \psi(s)$. These are compatibility-based
facts. A model giving $\varphi$ probability $0.99$, or even probability $1$
after assigning zero mass to compatible exceptions, is a different claim.

As observations accumulate, player histories extend by concatenation. The
compatible sets of complete histories at two different moments are not
generally nested: their members refer to different endpoints. Likewise, sets
of compatible current game states need not shrink while the game changes.
Each transition extends the actual history and appends its disclosed observations.
To compare possibilities before and after that transition, align their time or
decision endpoints; see \cref{sec:belief}.

\subsection{Why the full epistemic state is sufficient for continuation}

Although two histories can share $\epss(h)$, their discarded details do not
affect future behavior \emph{under this model}. Define
\begin{equation}
\mathsf T((s,i_0,i_1),a)
=\bigl(\nexts(s,a),\ i_0\cat O_0(s,a),\ i_1\cat O_1(s,a)\bigr).
\label{eq:epistemic-transition}
\end{equation}

\begin{proposition}[Continuation sufficiency]
Fix policies whose action distributions depend only on player information
states. Histories with the same full epistemic state induce the same law of
future full epistemic states and, when defined, terminal payoff.
\label{prop:sufficiency}
\end{proposition}
\begin{proof}
The common game state gives the same actor and legal actions. The actor's
common information state gives the same action distribution. For each realized
action, \cref{eq:epistemic-transition} gives the same successor tuple.
Induct on the number of future transitions. Terminal payoff depends on the
terminal game state, so its law agrees as well.
\end{proof}

The tuple therefore provides a lossless continuation representation for this
fixed-policy model. An implementation can use this result by retaining all
state context that affects transitions and disclosures, together with each
player's observation history. Histories remain a convenient foundation for
defining how those components arise.

\Needspace{18\baselineskip}\part{Probability and decision targets}

\section{Policies and probabilities of play}
\label{sec:policy}

The observation history tells us which actions a player can choose. To describe
how play proceeds, we also need a rule for choosing among them. After $\ell$,
one guesser policy always chooses orange; another chooses orange with
probability $3/4$ and blue with probability $1/4$.

Write $\Info_q^{\mathrm{act}}=\{i_q:D_q(i_q)\ne\bot\}$ and, on that set,
$A_q(i_q)=D_q(i_q)$. For a finite or countable set $X$, $\Delta(X)$ is the set of
probability distributions on $X$: nonnegative masses summing to one.
Action sets themselves remain finite.

\Needspace{7\baselineskip}
\begin{definition}[Behavioral policy]
A behavioral policy for player $q$ assigns a distribution over legal choices:
\[
\pi_q(\cdot\mid i_q)\in\Delta(A_q(i_q))
\quad\text{for every }i_q\in\Info_q^{\mathrm{act}}.
\]
\end{definition}

A pair $\pi=(\pi_0,\pi_1)$ is a \emph{policy profile}. At each decision, the
acting player's policy samples an action and deterministic advancement gives
the next state. The two guesser policies above therefore describe different
patterns of play, but each uses the same choice probabilities under orange
and blue whenever the guesser's observations agree.

\textit{Randomization.} Conditional on the full history, the next action has
the law $\pi_q(\cdot\mid I_q(h))$. Independent fresh draws realize this law;
persistent random memory or shared randomness needs a richer model if it
changes that conditional distribution.

\subsection{Path probability}

Fix an initial-state distribution $\lambda\in\Delta(\Init)$. A point mass
describes a fixed starting state. For a valid history $h_n$, its prefix
probability under $\lambda$ and $\pi$ is
\begin{equation}
\Prob_{\lambda,\pi}([h_n])
=\lambda(s_0)\prod_{j=0}^{n-1}
\pi_{p(s_j)}(a_j\mid I_{p(s_j)}(h_j)).
\label{eq:path}
\end{equation}
The notation $[h_n]$ means the event that play begins with this prefix.
There is no chance-transition factor. Histories at a fixed depth are disjoint
prefix events, but a history and its extension are not: $[h']\subseteq[h]$
when $h\preceq h'$. Therefore \cref{eq:path} is not a normalized mass function
over \emph{all} finite histories taken together.

\section{Weighting compatible histories}
\label{sec:belief}

\Needspace{7\baselineskip}
\begin{definition}[Belief at a player information state]
A belief at $i_q$ is a probability distribution
\[
b_q(\cdot\mid i_q)\in\Delta(\mathcal C_q(i_q)).
\]
It assigns zero mass outside compatibility. It need not assign positive mass
to every compatible history.
\end{definition}

A prior and a model of past play provide one way to choose the weights. In the
signal game, we ask how frequently each type of sender would choose $\ell$.

\begin{example}[A signal changes probability without eliminating possibility]
A sender privately knows a type $\otype$ or $\btype$. The observer's initial prior is
$1/2$ each. Both types can legally choose public signal $\ell$. Under a particular
sender policy,
\[
\Prob(\ell\mid \otype)=4/5,\qquad \Prob(\ell\mid \btype)=1/5.
\]
After $\ell$, both types remain compatible, but
\[
\Prob(\otype\mid \ell)=\frac{(1/2)(4/5)}{(1/2)(4/5)+(1/2)(1/5)}=4/5.
\]
\label{ex:signal-posterior}
\end{example}

\begin{principle}{From possibility to probability}
Observing $\ell$ leaves both orange and blue compatible. A sender model that chooses
$\ell$ more often when orange makes orange more probable. The observation determines
which histories fit; the model determines their relative weights.
\end{principle}

\subsection{The natural endpoint: encountering a decision}

\begin{lemma}[Decision-state compatibility is prefix-free]
\label{lem:prefix-free}
For $i_q\in\Info_q^{\mathrm{act}}$, no member of $\mathcal C_q(i_q)$ is a
proper prefix of another member.
\end{lemma}
\begin{proof}
Suppose $h\prec h'$ are both members. Decision consistency says that $q$
must act at the end of $h$. The extension $h'$ therefore contains an
additional action of $q$, so $\own_q(h')$ is strictly longer than
$\own_q(h)$. Own-choice recoverability says that equal information states
have equal own-action sequences, giving a contradiction.
\end{proof}

Thus a play can encounter a particular acting information state at most once.
Its compatible prefix events are disjoint. Let
\[
z_q(i_q)=\sum_{h\in\mathcal C_q(i_q)}\Prob_{\lambda,\pi}([h]).
\]
This is the probability that play encounters that decision information state.
When it is positive, the reach-conditioned posterior is
\begin{equation}
b_q^{\mathrm{reach}}(h\mid i_q)
=\frac{\Prob_{\lambda,\pi}([h])}{z_q(i_q)},
\qquad h\in\mathcal C_q(i_q).
\label{eq:reach-posterior}
\end{equation}
The denominator is at most one, even if compatible histories have different
lengths. This is the usual planner-root case: condition on arriving at the
decision the player is actually facing.

\subsection{Other history endpoints}

To derive a belief by conditioning generated play, first specify a random
finite history $H$. For example, stop at player $q$'s $k$th decision, provided
that decision occurs; the experiment can explicitly condition on its occurrence.
Alternatively choose a declared prefix-free frontier: a set of histories no
one of which is a proper prefix of another. Restrict and normalize the path
probabilities on that frontier when its total probability is positive.

Then, when $\Prob(I_q(H)=i_q)>0$,
\begin{equation}
b_q(h\mid i_q)=
\frac{\Prob(H=h)\ind\{I_q(h)=i_q\}}
{\sum_{g:\ I_q(g)=i_q}\Prob(H=g)}.
\label{eq:posterior}
\end{equation}
The law of $H$ includes the prior, policy profile used to generate the past,
and endpoint selection rule. \Cref{lem:prefix-free} supplies a canonical
reach event for acting information states. For a waiting information state,
unobserved transitions can leave $I_q$ unchanged, so compatible prefixes can
overlap; the general experiment must specify which endpoint is sampled.

An endpoint is a modeling device, not necessarily a timestamp disclosed to the
player. Conditioning on an undisclosed global step count would change their
information if the model treated that count as known. A decision occurrence
recoverable from the player's own action chronology is often a cleaner
endpoint. At a probability-zero observation, ordinary conditioning does not
specify a posterior; an explicit off-path model is needed.

\subsection{Derived beliefs preserve joint relationships}

Several compatible histories may end at the same current game state. To find
that state's probability, add the probabilities of those histories. More
generally, mapping histories through a function $f$ produces a distribution
called the \emph{pushforward}, $f_{\#}b$:
\[
(f_{\#}b)(x)=\sum_{h:f(h)=x}b(h).
\]
Thus $\last_{\#}b_q$ is a belief over current game states and
$\epss_{\#}b_q$ is a belief over full epistemic states. The latter keeps the
correlation between game configuration and opponent information. Under
\cref{prop:sufficiency}, it is sufficient for continuation under the fixed
information-state policies. The game-state marginal alone need not be.

\subsection{Own choices versus observed opponent choices}

At a decision information state $i_q$, the player's own action probability
$\pi_q(a\mid i_q)$ is the same throughout $\mathcal C_q(i_q)$. If it is
positive, conditioning solely on the realization of that own action does not
reweight those hypotheses: the common likelihood factor cancels. Subsequent
disclosures can of course reweight them.

An opponent action is different. Compatible histories can give that opponent
different information states and hence different action probabilities. The
observed action can therefore change the player's posterior through the
opponent model. This is inference from behavior, not an engine disclosure of
the opponent's hidden information.

\subsection{Advancing and updating a belief}
\label{sec:belief-update}

Add a preliminary decision to the signal game: the guesser can request a
signal or leave for a draw. Before choosing, it assigns mass $1/2$ to each
type. Fix its choice to \emph{request}. That advances both hypotheses to the
sender's decision, preserving their weights. The sender then chooses a signal
with the policy from \cref{ex:signal-posterior}.

\begin{center}
\begin{tabular}{@{}ccccc@{}}
\toprule
Type & Mass after request & Sender choice & Likelihood & Branch mass\\
\midrule
$\otype$ & $1/2$ & $\ell$   & $4/5$ & $2/5$\\
$\otype$ & $1/2$ & $\mathsf r$ & $1/5$ & $1/10$\\
$\btype$ & $1/2$ & $\ell$   & $1/5$ & $1/10$\\
$\btype$ & $1/2$ & $\mathsf r$ & $4/5$ & $2/5$\\
\bottomrule
\end{tabular}
\end{center}

Each branch contains an advanced game state and both updated observation
histories. On observing $\ell$ and receiving the guess request, the guesser keeps
the two $\ell$ branches. Their masses sum to $1/2$; normalizing gives $4/5$ for
$(s^{\otype},\ell,t^{\otype,\ell})$ and $1/5$ for $(s^{\btype},\ell,t^{\btype,\ell})$, with the preliminary
request prefix understood.

The update has three operations: advance each hypothesis under the chosen
action, branch and weight by intervening opponent behavior, then match the
new observation and normalize. The first operation changes the possible
worlds; the latter two change how much probability each surviving world carries.
The next guess uses this new distribution at the new decision state.

\section{Payoff and target values}
\label{sec:value}

\Needspace{7\baselineskip}
\begin{definition}[Terminal payoff]
For $q\in P$, let $u_q:F\to\{-1,0,+1\}$ assign loss, draw, and win,
respectively. We assume zero-sum payoff: $u_0(s)=-u_1(s)$.
\end{definition}

The signal game's guesser receives $+1$ for a correct guess and $-1$ for
an incorrect guess; the sender receives the opposite payoff. Expected payoff
therefore measures the difference between win and loss probabilities.

Define $G_q^\pi(h)$ whenever continuation from $h$ under $\pi$ reaches a
terminal state almost surely. It is the expected terminal payoff to $q$.
At terminal $h$, it equals $u_q(\last(h))$. At nonterminal $h$, write
$s=\last(h)$ and $r=p(s)$. Conditioning on the next action gives
\begin{equation}
G_q^\pi(h)=\sum_{\substack{a\in A(s)\\\pi_r(a\mid I_r(h))>0}}
\pi_r(a\mid I_r(h))\,
G_q^\pi\bigl(h\cat(a,\nexts(s,a))\bigr).
\label{eq:history-value}
\end{equation}
Positive-probability continuations inherit almost-sure termination. Restricting
the sum to them keeps every continuation value in the expression defined.
Bounded payoff makes this expectation finite even when expected game length
is infinite. A model admitting nonterminating play needs an additional payoff
convention for that play before assigning a terminal-value target to it.

\subsection{Belief value and action value}

The \emph{support} of a belief contains its positive-mass histories:
$\supp(b_q)=\{h:b_q(h)>0\}$. A compatible history can fall outside it. For
example, if $\ell$ has positive likelihood under orange and zero likelihood
under blue, blue remains rule-compatible after that signal but receives
posterior mass zero.
For a belief $b_q$ at $i_q$, define
\begin{align}
V_q(b_q,\pi)&=\sum_{h\in\supp(b_q)}b_q(h)G_q^\pi(h),
\label{eq:V}\\
Q_q(b_q,a;\pi)&=\sum_{h\in\supp(b_q)}b_q(h)
G_q^\pi\bigl(h\cat(a,\nexts(\last(h),a))\bigr).
\label{eq:Q}
\end{align}
For $V$, continuation must terminate almost surely from every history in this
support. For $Q$, $i_q\in\Info_q^{\mathrm{act}}$, $a\in A_q(i_q)$, and
continuation after that forced first action must terminate almost surely on
the support. Decision consistency makes $a$ legal in every summand.

$Q$ asks what happens when the player \emph{chooses} $a$ and then follows
$\pi$. It remains meaningful when $\pi_q(a\mid i_q)=0$: the first choice is
an intervention, while later choices follow the policy. The historical policy
used to infer $b_q$ can differ from the continuation policy being evaluated.
This lets a player use a model of past behavior to compare new plans.

\begin{proposition}[Value identities]
Let $i_q$ be an acting information state and $b_q$ a belief at $i_q$.
Suppose continuation under $\pi$ terminates almost surely on $\supp(b_q)$,
including after each forced legal action $a\in A_q(i_q)$. Then
\begin{align}
-1&\le V_q(b_q,\pi)\le1,\qquad -1\le Q_q(b_q,a;\pi)\le1,\\
V_q(b_q,\pi)&=\sum_{a\in A_q(i_q)}\pi_q(a\mid i_q)Q_q(b_q,a;\pi).
\label{eq:mixture}
\end{align}
\end{proposition}
\begin{proof}
The bounds follow by averaging payoffs in $[-1,1]$. Substitute
\cref{eq:history-value} into \cref{eq:V}. Each next-action probability is
constant on the compatibility class, so the finite action sum factors outside
the history sum. Under the stronger forced-action assumption, the zero-weight
action terms are also defined and can be included in \cref{eq:mixture}.
\end{proof}

\Needspace{11\baselineskip}
\begin{example}[Running game: choosing and mixing guesses]
At $i_\ell$, the posterior assigns probabilities $4/5$ to $\otype$ and $1/5$ to $\btype$.
Both guesses terminate immediately, so the action values are
\begin{align*}
Q_1(b_\ell,g_{\otype};\pi)&=\tfrac45(+1)+\tfrac15(-1)=\tfrac35,\\
Q_1(b_\ell,g_{\btype};\pi)&=\tfrac45(-1)+\tfrac15(+1)=-\tfrac35.
\end{align*}
If the guesser chooses $g_{\otype}$ with probability $3/4$, its policy value is
$V_1(b_\ell,\pi)=\tfrac34\tfrac35+\tfrac14(-\tfrac35)=\tfrac3{10}$.
Thus $3/5$ is the value of choosing orange, while $3/10$ is the value of this
randomized continuation policy at the same information state.
\end{example}

\subsection{Defining the value target}

To assign an expected payoff to an information state, fix a belief over
compatible histories and policies for subsequent play. These choices specify
the distribution of terminal outcomes. Writing the belief construction as
$B_q:i_q\mapsto b_q$, the resulting target is
$i_q\mapsto V_q(B_q(i_q),\pi)$. Changing $B_q$ or $\pi$ changes the question.
A best response to a fixed opponent, a worst-case criterion, and an equilibrium
objective each add their own optimization problem to this basic expectation.

Under one common distribution of terminal play, zero-sum payoffs give opposite
expectations. Players with different observations may instead use different
subjective distributions. Both can then assign themselves positive expected
value; each expectation describes that player's model of the future.

\section{Evaluators and representations of their inputs}
\label{sec:evaluator}

Search can use an evaluator to score a position without continuing play to
termination. We begin with a scorer that receives the player's observation history.

\Needspace{7\baselineskip}
\begin{definition}[Information-state evaluator]
An information-state evaluator for player $q$ is a scoring function
\[
e_q:\Info_q\to\mathbb R,
\]
or a function restricted to a specified set of information states, such as
nonterminal decision contexts.
\end{definition}

The project also calls this a \emph{state evaluator}. The scoring function can
be hand-designed, learned, or composite. To train one for the
randomized policy in the running example, the exact target at $i_\ell$ is $3/10$.
To score the always-orange policy at the same input, the target is $3/5$.
The belief and continuation policy identify which of these tasks is intended.

A \emph{sampled-world evaluator} has a game-state input, for example
$g_q:S\to\mathbb R$ on a specified nonterminal domain. An evaluator on the full
epistemic tuple can additionally use both players' represented histories.
These richer inputs are useful inside trusted simulation; an information-state
evaluator receives only the selected player's observation history.

\begin{principle}{Evaluating a hypothetical player's information}
For a sampled history $h$, an information-state evaluator is called on $I_q(h)$.
Projection extracts that hypothetical player's experience before scoring it.
The scorer therefore uses the same input type in real and hypothetical play.
\end{principle}

The implementation encodes this input before passing it to the scorer.
\Cref{sec:representations} develops the representation map and shows when a
smaller encoding retains everything needed for a particular target.

\subsection{Range and target identity}

An evaluator's range is $\mathbb R$ unless restricted. A search that accepts
only $[-1,1]$ settlements can require this range or apply a declared
transformation. Clipping and $\tanh$ produce new scoring functions; for example,
rescaling a tactical heuristic into $[-1,1]$ supplies a bounded search score,
while estimating terminal payoff still requires a specified target and validation.
\Cref{sec:cutoff} works through a root-information evaluator called at an
opponent's decision.

\Needspace{18\baselineskip}\part{Simulation and search}

\section{Hypothetical simulation}
\label{sec:simulation}

At an actual root decision, fix root player $q$, information state $i_q$, and
belief $b_q$. A simulation first samples $h\sim b_q$, choosing one possible
history to explore under the root's model.
In a computational approximation, a finite weighted collection of hypotheses,
often called \emph{particles}, supplies this sampling distribution.
An exact full epistemic representation can replace the history for continuation
under \cref{prop:sufficiency}; a game-state sample alone is insufficient when
opponent memory matters.

At a nonterminal simulated history, $r=p(\last(h))$ acts using
\[
a\sim\mu_r(\cdot\mid I_r(h)),\qquad
h\leftarrow h\cat(a,\nexts(\last(h),a)).
\]
This equation describes simulation under a genuine continuation policy profile.
The sampled opponent acts on its own hypothetical observation history.
Search will use the returns from these simulations to choose which actions
to explore. One question will then arise: the simulated opponent acts on its
own information, while the search record is built from the root's information.
\Cref{sec:opponent-model} examines that relationship after we have seen the
record and selector at work.

\subsection{Hypothesis consistency and failure}

Sample the game configuration and both player histories jointly so that one
compatible history witnesses the tuple. After each action,
both observation histories are updated with the same transition's respective
observation sequences.

In the theoretical game, a legal action advances successfully. An adapter may
fail to represent, bind, or execute it. Such a failure is an implementation
outcome, not a new legal game outcome. Rejection, timeout, unsupported
representation, and stopped execution must remain separately typed.

\subsection{Cutoffs and settlement}

A simulation can stop at terminal play, at a declared computational cutoff,
or at an execution failure. Terminal play supplies $u_q(s)$. A nonterminal
cutoff requires a declared scoring or continuation rule. A failure supplies no
strategic value.

A \emph{rollout} continues from a search leaf using declared action-selection
policies. It uses the same game and observation rules. Its endpoint can be
terminal, another cutoff, or failure. ``Leaf'' means that represented graph
exploration stops for this iteration, not that the game has ended.

Record a successful settlement as $(v,\kappa)$, with origin $\kappa$ identifying
terminal payoff, a specific evaluator, or a specified bounded-rollout rule.
The computation horizon, continuation policies, evaluator configuration, and
payoff perspective contribute to its meaning. For example, a terminal $+1$ records a win in the hypothetical game; an
evaluator score of $+1$ records that evaluator's judgment at its cutoff.

\section{Search nodes, action branches, and statistics}
\label{sec:graph}

A simulation follows a concrete hypothesis. A search structure determines which
hypotheses contribute to the same decision statistics.

\Needspace{7\baselineskip}
\begin{definition}[Acting-player information-state node]
For a nonterminal history $h$, let
\[
K(h)=\bigl(p(\last(h)),I_{p(\last(h))}(h)\bigr).
\]
A node $n=(r,i_r)$ groups all histories with the same acting player and the
same information state for that player. This group is a \emph{fiber} of $K$.
Its legal action set is $A(n)=A_r(i_r)$.
\end{definition}

Decision consistency makes the action set well-defined. Terminal histories
produce outcome records rather than decision nodes. An action branch is the
pair $(n,a)$, with $a\in A(n)$. It may have several successor decision nodes
and terminal outcomes, since different histories in the fiber can produce
different consequences for the same action.

\subsection{Search-local statistics}

The search record $Z$ stores the actions explored and their returns. For each
node-action pair, retain
\begin{align*}
N(n,a)&=\text{number of value-bearing traversals},\\
W(n,a)&=\text{sum of their backed root-perspective values}.
\end{align*}
When $N(n,a)>0$, set $\widehat Q(n,a)=W(n,a)/N(n,a)$.
For an unvisited action the mean is undefined; selection uses its declared
exploration rule. Attempt, failure, proposal, and origin counts describe the
additional computation that produced these value-bearing traversals.

Each search invocation starts with a fresh record.

\subsection{Watching the guesser search}
\label{sec:example}

Fix root $q=1$, belief $b_\ell(\otype)=4/5$, $b_\ell(\btype)=1/5$, terminal settlement, and
one-step exploration. Try unvisited actions in order $g_{\otype},g_{\btype}$, then choose
the highest current mean, breaking ties toward $g_{\otype}$. The following is one possible sample
sequence under $b_\ell$. A cell records $(N,W,\widehat Q)$ \emph{after} the attempt.

\begin{center}
\small
\begin{tabular}{@{}clclll@{}}
\toprule
Attempt & Type & Choice & Settlement & $g_{\otype}$ statistics & $g_{\btype}$ statistics\\
\midrule
1 & $\otype$ & $g_{\otype}$ & $+1$, terminal & $(1,1,1)$ & $(0,0,\text{--})$\\
2 & $\otype$ & $g_{\btype}$ & $-1$, terminal & $(1,1,1)$ & $(1,-1,-1)$\\
3 & $\btype$ & $g_{\otype}$ & $-1$, terminal & $(2,0,0)$ & $(1,-1,-1)$\\
4 & $\otype$ & $g_{\otype}$ & $+1$, terminal & $(3,1,1/3)$ & $(1,-1,-1)$\\
5 & $\otype$ & $g_{\otype}$ & $+1$, terminal & $(4,2,1/2)$ & $(1,-1,-1)$\\
6 & $\btype$ & $g_{\otype}$ & Execution failure & $(4,2,1/2)$ & $(1,-1,-1)$\\
\bottomrule
\end{tabular}
\end{center}

After attempt 2, $1>-1$, so the selector chooses orange. After attempt 3 its
mean falls to zero, still above blue's $-1$. Six attempts have consumed the
budget, and five have supplied values. The highest-mean final rule chooses
$g_{\otype}$. Its backed mean $1/2$ differs from the exact target $3/5$ because this
short realized sample contains only four orange-action returns.

The sixth attempt produces a typed failure and leaves both value counts
unchanged. In the toy mathematics, we can calculate the unexecuted continuation:
guessing orange under $\btype$ would pay $-1$. Had execution completed, orange's statistics
would become $(5,1,1/5)$. This comparison explains the missing contribution;
the implementation records the observed failure, without substituting that
counterfactual payoff. A single trace illustrates arithmetic. The systematic
selective-failure model in \cref{sec:estimates} explains bias across repetitions.

\subsection*{Connecting decisions in longer games}

The guesser's choice ends the signal game. In a longer game, an action can
lead to another decision node, and different hypotheses may produce different
successors. We therefore need to describe how these nodes connect.

\Needspace{7\baselineskip}
\begin{definition}[Search connectivity]
A successor node $n'$ is possible from $(n,a)$ if there exists a valid history
$h$ with $K(h)=n$ such that, writing
$h'=h\cat(a,\nexts(\last(h),a))$, the successor is nonterminal and $K(h')=n'$.
\end{definition}

Each edge has a history that produces it, but a path through the graph may
combine edges from different histories. Simulation follows the successor of
its current hypothesis, which keeps the path consistent with one possible game.

\begin{example}[Why the quotient need not be a tree]
Consider two initial states with a hidden order of play. Each player makes one
private choice; the other receives no observation of its content. In one
initial state player $0$ acts first; in the other player $1$ acts first. Each
player receives the same first-decision request when their turn arrives.
The first-decision nodes can have an edge $n_0\to n_1$ in the first case and
$n_1\to n_0$ in the second. The quotient graph has a cycle, although each
actual history ends after two choices. Traversing that apparent cycle would
splice incompatible order hypotheses. Own-choice recall prevents an actual
player's second choice from being confused with their first.
\end{example}

This example concerns the graph across the allowed history family; a particular
root's support can exclude some edges. Unfolding valid histories gives a tree
of occurrences. Sharing information-state keys gives a graph. Neither data
structure removes the need to retain a consistent simulated world.

\subsection{How a node acquires a population}

The node key specifies an information state, not a distribution over its
compatible histories. A root belief, earlier selection policies, observed
continuations, and sampling decisions determine the distribution of histories
that reach a node during a particular search.

If $\nu$ is a declared distribution over histories in node $n$, its successor
distribution after fixing $a$ is obtained by pushing $\nu$ through advancement
and $K$, with terminal outcomes represented separately. Uniformly choosing a
successor node is not implied by the connectivity relation. Distinct nodes can
aggregate different amounts of probability.

\section{A complete search iteration}
\label{sec:algorithm}

A search fixes its root belief and configuration and uses root-perspective
payoffs. Algorithm~\ref{alg:search} makes $B$ attempts, calling
Algorithm~\ref{alg:attempt} to explore at most $D$ actions before settlement.
Fallible operations return $\textsc{Ok}(x)$ or
$\textsc{Failure}(\text{reason})$.

\begin{algorithm}[H]
\caption{Bounded search with short-circuit failure handling}
\label{alg:search}
\small
\begin{algorithmic}[1]
\Require Root $(q,i_q)$, belief $b_q$, configuration $\xi$, finite bounds $B,D$
\State Initialize search record $Z$; on failure, return declared legal fallback
\For{$j=1,\ldots,B$} \Comment{Each entry consumes one attempt, including failures}
  \State $\eta\gets\textsc{Attempt}(q,i_q,b_q,Z,\xi,D)$
  \If{$\eta=\textsc{Ok}(L,v,\kappa)$}
    \State Back up $v$ along $L$; retain settlement origin $\kappa$
  \Else
    \State Record the failure reason; perform no numeric backup
  \EndIf
\EndFor
\State \Return Root distribution $\rho(\cdot\mid i_q,Z)$, with legal fallback
\end{algorithmic}
\end{algorithm}

Inside \textsc{Attempt}, the notation
$x\gets\textsc{Try}(f)$ evaluates $f$, extracts $x$ only from
$\textsc{Ok}(x)$, and otherwise \emph{immediately returns} that failure from
the attempt. This short-circuit rule applies before any later line runs.
All operations, including initialization, sampling, final choice, and settlement,
have finite work bounds and typed failure or a declared fallback.

\begin{algorithm}[H]
\caption{One search attempt}
\label{alg:attempt}
\small
\begin{algorithmic}[1]
\Function{Attempt}{$q,i_q,b_q,Z,\xi,D$}
  \State $h\gets\textsc{Try}(\operatorname{Sample}(b_q))$; $L\gets[\,]$
  \For{$d=0,\ldots,D$}
    \If{$\last(h)\in F$}
      \State \Return $\textsc{Ok}(L,u_q(\last(h)),\textsc{Terminal})$
    \EndIf
    \If{$d=D$}
      \State $(v,\kappa)\gets\textsc{Try}(\operatorname{Settle}(h,q,\xi))$
      \State \Return $\textsc{Ok}(L,v,\kappa)$
    \EndIf
    \State $n\gets\textsc{Try}(\operatorname{ProjectNode}(h))$
    \State $M\gets\textsc{Try}(\operatorname{Propose}(n,Z,\xi))$
    \State $(a,\mathit{new})\gets\textsc{Try}(\operatorname{Select}(n,M,Z,\xi))$
    \State $h'\gets\textsc{Try}(\operatorname{Advance}(h,a))$
    \State Append $(n,a)$ to $L$; set $h\gets h'$
    \If{$\mathit{new}$}
      \State $\textsc{Try}(\operatorname{RegisterSuccessor}(n,a,h,Z))$
      \State $(v,\kappa)\gets\textsc{Try}(\operatorname{Settle}(h,q,\xi))$
      \State \Return $\textsc{Ok}(L,v,\kappa)$
    \EndIf
  \EndFor
\EndFunction
\end{algorithmic}
\end{algorithm}

\code{ProjectNode} constructs $K(h)$ in the supported representation;
\code{Propose} and \code{Select} use the node's player information and the
search record. Advancement and successor projection operate on the trusted
hypothesis. \code{Settle} checks terminal payoff first, then runs the declared
evaluator or bounded continuation; failure propagates through the same
\textsc{Try} rule. Structural changes made before a failed attempt may remain
in $Z$ as recorded exploration, while value counts increase only at backup.
A cutoff at depth zero can settle the root without evaluating any root action,
so the final choice rule also handles an empty value-bearing path.

\subsection{Selection from the current search record}
\label{sec:selector-policy}

At a node $n=(r,i_r)$, a computational selector returns a distribution over
represented actions:
\[
\alpha_r(\cdot\mid i_r,Z,\xi)\in\Delta(A_{\mathrm{rep}}(n)).
\]
Holding search record $Z$ and configuration $\xi$ fixed, it is a function of
the node, not of inaccessible details of the current hypothesis.
Hypothesis-dependent menu filtering or choice features can violate this
fixed-record requirement even when the node keys themselves are safe.

\Needspace{13\baselineskip}
An illustrative UCB-style score for visited actions is
\begin{equation}
\epsilon_q(n)\widehat Q(n,a)
+c\sqrt{\frac{\log(1+N(n))}{N(n,a)}},\qquad
N(n)=\sum_aN(n,a),
\label{eq:ucb}
\end{equation}
where $c\ge0$ and $\epsilon_q(n)$ is $+1$ for a root-player node and $-1$
for an opponent node. The first term favors larger observed returns at root-player
nodes and smaller ones at opponent nodes. For $c>0$, the second gives
less-sampled actions a larger exploration bonus, encouraging the search to
revisit alternatives to an early favorite. Increasing $c$ places more weight
on exploration.

Unvisited actions need an initial exploration rule, and equal scores need a
tie-break. This score illustrates the bandit-guided allocation used in
UCT~\cite{kocsis2006}; \cref{sec:related} discusses the conditions of its analysis.

As $Z$ evolves, the selection distribution changes with the accumulated
statistics. At opponent nodes, the minimizing sign defines the search-conditioned
response described in \cref{sec:opponent-model}. Prediction under a fixed
behavioral opponent instead uses its policy. The model used to interpret past
opponent choices may differ from the rules used to choose future actions
during search and rollout.

\subsection{Which opponent is being simulated?}
\label{sec:opponent-model}

A behavioral opponent policy $\mu_r(\cdot\mid i_r)$ gives the same action
distribution at equal opponent information. The selector instead conditions
on the search record as well: $\alpha_r(\cdot\mid i_r,Z,\xi)$.
Across root searches, $Z$ can carry private information that changes the
response at the same opponent node.

\begin{proposition}[Local selector invariance permits an opponent information advantage]
\label{prop:root-conditioned}
Sharing opponent information-state keys and obeying fixed-record selector
invariance need not produce one behavioral opponent policy across root contexts.
\end{proposition}
\begin{proof}[Counterexample]
Run two sender-root searches in the signal game: the sender knows type $\otype$
in one and type $\btype$ in the other. In each search, consider sending $\ell$.
The guesser then faces the same information state $i_\ell$. Each root search
evaluates both guesses against its known type and stores their sender payoffs
in $Z_{\otype}$ or $Z_{\btype}$. Minimizing sender payoff chooses $g_{\otype}$ under $Z_{\otype}$ and
$g_{\btype}$ under $Z_{\btype}$. The selector uses only $(i_\ell,Z)$, satisfying fixed-record
invariance. A single guesser policy $\mu_1(\cdot\mid i_\ell)$ would have to give
both guesses probability one, which is impossible.
\end{proof}

\begin{principle}{Two information boundaries, two useful questions}
Root perspective safety asks whether the final choice depends only on the
root's legitimate information. Opponent realizability asks whether the modeled
opponent's choices agree across all contexts with the same opponent information.
The counterexample satisfies the first condition while failing the second.
\end{principle}

For prediction under a specified opponent policy, simulation can use that
policy directly. An optimization procedure seeking a consistent strategic
opponent must constrain the induced kernels across the relevant information
sets and root contexts. Local maximization or minimization of search statistics
alone supplies a search-conditioned response. Cowling, Powley, and Whitehouse
discuss the closely related private-card information advantage when the
searching player's cards remain fixed across determinizations
\cite[Section~X, pp.~140--141]{cowling2012}.

\subsection{Expansion and action coverage}

Write $A_{\mathrm{rep}}(n)\subseteq A(n)$ for represented actions.
Expansion may introduce all actions or grow this set progressively. A proposal
rule is computational machinery; an omitted action remains legal. The actual
simulation must bind a proposed semantic action to the corresponding action in
its current hypothesis without changing its meaning.

This separates four implementation questions: is an action legal, was it
proposed, was it admitted to the configured search, and did its attempted
execution succeed? Their sets can differ without changing the theoretical
legal set. When reporting a result, identify which legal actions the search
considered.

\subsection{Backup and root choice}

For each occurrence $(n,a)$ on the exploration path receiving settlement $v$,
\[
N(n,a)\leftarrow N(n,a)+1,\qquad W(n,a)\leftarrow W(n,a)+v.
\]
The same root-perspective scalar is stored at every node. Opponent selection
applies its sign when reading that scalar, as in \cref{eq:ucb}. An acting-player
storage convention is also possible if storage and selection use matching
conversions.

Rollout choices need not create graph visits; their role is to supply the
settlement for the recorded exploration path. Any scheme that also updates
rollout nodes is an additional algorithm choice.

After budget $B$, let $Z_B$ be the random search record. A root choice rule
returns $\rho(\cdot\mid i_q,Z_B)\in\Delta(A_q(i_q))$. Examples include highest
mean, highest visit count, and a declared distribution based on statistics.
Ties, unvisited actions, and zero successful iterations require explicit
handling. A legal fallback makes the procedure a total policy; without one,
a procedure that can fail is only a partial policy implementation.

The induced search policy is
\begin{equation}
\pi_q^{\mathrm{search}}(a\mid i_q)
=\Expect[\rho(a\mid i_q,Z_B)],
\label{eq:search-policy}
\end{equation}
where expectation includes sampling, exploration, and any final choice
randomness. All persistent configuration and the budget convention are fixed
when reading this expression. \Cref{prop:search-safety} supplies sufficient
conditions for this expression to depend only on $i_q$.

\subsection{Evaluating a cutoff at the opponent's turn}
\label{sec:cutoff}

Now start a different search at the sender's decision, with root $q=0$ and
known type $\otype$. Consider the branch sending $\ell$ and stop at the next state,
where the guesser acts. Fix the continuation model: the guesser always
chooses $g_{\otype}$. Define the sender's information-state evaluator on these waiting information
states by its expected terminal payoff under that continuation:
\[
e_0(I_0(h_{\otype,\ell}))=-1,\qquad e_0(I_0(h_{\btype,\ell}))=+1.
\]
The sender's own history identifies its type and sent signal. Any compatible
history on this domain has the same terminal payoff under this policy,
so every belief supported there yields the displayed value. At the $\otype$ cutoff,
\code{Settle} projects $I_0(h_{\otype,\ell})$, calls $e_0$, and backs $-1$ to the
sender's $\ell$ branch with origin \textsc{RootInformationEvaluator}. The value is
exact for this continuation, while its origin records that simulation
stopped before the guesser chose.

\Needspace{9\baselineskip}
The guesser's evaluator at $i_\ell$, using $b_\ell$ and the same always-orange policy,
would return $e_1(i_\ell)=3/5$. Negating it would give $-3/5$, which differs from
the sender's $-1$ because the guesser is uncertain about type. Calling the
root-information evaluator preserves the root's subjective target even when
the opponent is the next actor. A different settlement design may use a
sampled-world scorer or a common-model rollout, with its own target.

\section{Connecting search averages to a target}
\label{sec:estimates}

We can now relate the target, the evaluator, and the search average.
\begin{center}
\begin{tabularx}{\textwidth}{@{}lX@{}}
\toprule
Object & Meaning\\
\midrule
$Q_q(b_q,a;\pi)$ & A specified terminal-payoff expectation under a belief and continuation profile.\\
$e_q(i_q)$ & A scoring function on player information, with a separately specified target.\\
$\widehat Q(n,a)$ & The average of the values actually backed to a node-action pair.\\
\bottomrule
\end{tabularx}
\end{center}

Sampling, continuation, settlement, and the population reaching a node determine
when a backed mean estimates a specified value target. The following cases show
how to establish that connection and how it can change.

\subsection{A fixed-action reference experiment}

Suppose an experiment repeatedly samples $h\sim b_q$, fixes the same root
action $a$, then follows the same continuation profile $\pi$ to terminal payoff.
Assume independent repetitions, almost-sure termination, and a fixed,
deterministic number $m\ge1$ of repetitions. Each return has expectation
$Q_q(b_q,a;\pi)$, and their mean is an unbiased estimate of that target. Since returns are bounded, the sample mean converges almost surely
to that target as the number of independent repetitions grows.

This fixed-action evaluation gives a useful reference experiment. In adaptive
search, node visitation selects a population, deeper choices evolve, and stopping
or execution can select which returns appear in the final mean. Each mechanism
can be analyzed against the reference target.

\subsection{Visitation changes the population}

Consider a variant of the signal game with different signal probabilities.
Let a root belief assign hidden type $\otype$ and $\btype$ probability $1/2$ each.
Suppose a simulated opponent who knows its type chooses an observed branch
$\ell$ with probabilities $0.9$ and $0.1$, respectively. Conditional on reaching
that branch, the hidden-type proportions are $0.9$ and $0.1$.

If a later fixed choice pays the root $+1$ under $\otype$ and $-1$ under $\btype$, then
\[
\Expect[v]=0\quad\text{under the root type prior},\qquad
\Expect[v\mid \ell]=0.9-0.1=0.8.
\]
The second average describes the branch-conditioned population under that
behavior model. Labeling it with its conditioning event explains why it differs
from the prior expectation.

There is an equally important invalid version of the example. If one node
represents the \emph{same acting-player information state} for both types,
and its selection rule directly uses hidden type to choose $\ell$ at those rates,
the rule uses a distinction absent from the actor's information. Here the
branch proportions arise from an information leak in selection.

\subsection{Adaptive continuation, stopping, and execution}

Even with independent returns and a fixed target, a data-dependent stopping
rule can change the expected final mean. Let each return independently be
$+1$ or $-1$, equally likely. Stop after one sample when the first is $+1$;
otherwise take a second sample. The possible means are $1$, $0$, and $-1$,
with probabilities $1/2$, $1/4$, and $1/4$. Thus
\[
\Expect[\overline v_{\mathrm{stop}}]
=\tfrac12(1)+\tfrac14(0)+\tfrac14(-1)=\tfrac14,
\]
although every individual return has expectation zero. This effect comes from
selecting when to report the average.

If continuation policy $\pi^{(j)}$ changes across successful iterations, the
$j$th return can target a different conditional value. The pooled mean then
combines a sequence of continuation targets. A convergence analysis must account
for how that process evolves.

Similarly, suppose half the hypotheses would return $+1$ and half $-1$, but
all negative-return hypotheses fail before producing a value. Successful
backups average $+1$. No failure was numerically scored as a win, yet the
surviving population is biased for the intended all-hypothesis target.

\subsection{Budgets, reuse, and uncertainty}

A wall-clock budget can produce a different sample allocation from a fixed
simulation budget, especially when runtime depends on the hypothesis. Report
the actual budget and cutoff mechanisms so that work counts can be related to
the completed trajectories they produced.

Reusing statistics under a new root belief requires more than matching keys.
For example, returns collected under one distribution of remembered opponent
information cannot simply be relabeled as samples from a different one.
Importance weighting would require adequate support, known relevant sampling
laws, and stable target semantics; this paper does not prescribe such a repair.
The default formal search therefore starts with search-local statistics.

Confidence intervals also need a sampling model. Adaptive visits and selected
successes are not automatically eligible for independent-sample formulas.
The arithmetic establishes the empirical mean; the sampling model determines
which uncertainty statement accompanies it.

\Needspace{18\baselineskip}\part{Representations and implementation}

\section{Representation contracts}
\label{sec:representations}

The computer represents these mathematical objects with data structures. An
engine object can encode a game state; a snapshot, event ledger, and knowledge
summary can encode player information. An engine world together with both
player records can encode the full epistemic tuple.

\subsection{Perspective safety as factorization}

Suppose an implementation constructs a policy-facing representation
$\Phi_q:\Hist\to\mathcal R_q$. Perspective safety requires that it factor through
the player information state:
\begin{equation}
\Phi_q(h)=r_q(I_q(h))
\quad\text{for some }r_q.
\label{eq:safety}
\end{equation}
Equivalently, indistinguishable histories produce equal policy-facing
representations.
If representation construction randomizes, the corresponding requirement is
equality of output distributions conditional on equal information states.

A representation can add computed consequences of observations without breaking
this condition. But labels must still distinguish an entailed fact from a
model-dependent belief. Randomization, hashes, identifiers, action order, cached
annotations, and diagnostic fields are all potential channels and belong to
the same policy-facing boundary.

\subsection{The representation map}

An implementation encodes information through
\[
r_q:\Info_q\to\mathcal R_q
\]
and applies $\widehat e_q:\mathcal R_q\to\mathbb R$. Its induced theoretical
function is $\widehat e_q\circ r_q$. Hence
\[
r_q(i)=r_q(i')\Longrightarrow
\widehat e_q(r_q(i))=\widehat e_q(r_q(i')).
\]
No more powerful fitting algorithm can recover a distinction that the input
representation has erased. A sufficient representation for one particular
target can be incomplete for other tasks or for full information history.

\begin{example}[Safe but incomplete]
Two histories show the same current public board and hand sizes. One includes
an earlier reveal of Lightning Bolt; the other includes an earlier reveal of
Mountain. A representation retaining only the current board and counts can be
perspective-safe but merges these information states. A scorer on that
representation must give them the same score.
\end{example}

\subsection{Completeness and task sufficiency}

A representation is complete for theoretical player history when $r_q$ is
injective on realizable information states:
\[
r_q(i)=r_q(i')\Longrightarrow i=i'.
\]
Combined with perspective safety, equality of representations is exactly
equality of information states. Completeness concerns preservation of
distinctions, not a requirement to store a literal uncompressed transcript.
An exact reversible encoding qualifies.

Decision sufficiency is weaker: there exists a decoder $\widehat D_q$ with
$\widehat D_q(r_q(i))=D_q(i)$. Target sufficiency is also weaker: for a particular
target $f$, there exists $\widehat f$ with $\widehat f\circ r_q=f$. These
requirements allow a representation to omit parts of the history while
preserving everything needed for a particular decision or scoring task.

\begin{example}[Omitting an irrelevant observation]
Extend the signal game by disclosing a decorative token $d\in\{\circ,\square\}$
at setup. Its distribution is independent of type, and the fixed sender model
ignores it. At a guesser decision, let $r_1(i)$ retain only the observed signal,
decision status, and legal menu. Histories with different $d$ then have the same
encoding, so $r_1$ is noninjective. Yet the menu is decoded exactly, and the
posterior over type depends only on the signal.

For the always-orange continuation, the target $f(i)$ is $3/5$ after $\ell$ and
$-3/5$ after $\mathsf r$. Defining $\widehat f$ by these two values gives
$f=\widehat f\circ r_1$: the compression is decision-sufficient and
target-sufficient. Retaining the token would add no predictive information for
this task. A sender model that responds to the token would change that
conclusion, because signal likelihoods could then depend on the omitted detail.
\end{example}

\subsection{Execution fidelity}

Executing an action in the implementation should produce the encoded version
of the state prescribed by the game model. Let $c:S\to\mathcal X$ encode a
game state, and let $\operatorname{bind}(s,a)$ give the engine's executable
representation of action $a$ in state $s$. For supported inputs, we require
\[
\widehat{\nexts}(c(s),\operatorname{bind}(s,a))=c(\nexts(s,a)).
\]
A refusal is an explicit failure, not an arbitrary alternative successor.
The engine may need world-specific identifiers to execute an action. The
policy-facing action representation must still respect the player's
information boundary.

If the engine representation is richer than $c(s)$, the relation is better
written using an abstraction map from engine objects to theoretical states.
The requirement is still that represented legal execution commute with the
theoretical transition. Extra engine bookkeeping does not establish extra
player knowledge.

\subsection{Incremental history and knowledge summaries}

The stored player record can be updated as observations arrive. On reachable
inputs, updating that record with new observations $w$ should give the same
result as encoding the full extended history $i\cat w$:
\[
\widehat U_q(r_q(i),w)=r_q(i\cat w).
\]
If the representation omits details required by later updates, it may fail this
condition even when it is sufficient to score the current position.

A remembered-knowledge ledger efficiently summarizes consequences of player
history. For a known-card draw, the update reads the previously established
card identity, combines it with the new draw disclosure, and records the new
location. An explicit update order preserves this derivation. Labels for facts
and model-dependent inferences retain the distinction developed in
\cref{sec:compatibility}.

\subsection{Simulation and search-state isolation}

Forking a sampled world must preserve both players' valid pasts while allowing
independent future updates. Shared mutable histories, belief weights, or search
statistics can contaminate otherwise valid branches. A world snapshot alone
does not reconstruct either player's remembered experience.

An implementation key $\widehat K$ realizes the chosen theoretical identity
$K$ when its equality classes match those of $K$. A coarser key pools additional
information states. The sufficient-compression example shows why such pooling
can be useful for a particular task; its adequacy depends on which decisions and
continuation targets the abstraction preserves.

A finite digest alone cannot establish injectivity over an unbounded history
space. An implementation that uses digests for keys therefore needs a way to
handle collisions and determine semantic equality.

\subsection{Composing a perspective-safe root policy}

The root receives legitimate information and constructs its own hypothetical
worlds. Those worlds can contain hidden details because they were generated by
the root's model. A complete safety argument follows their provenance through
the computation, including stopping and fallback behavior.

\begin{proposition}[End-to-end root perspective safety]
\label{prop:search-safety}
Consider a search that terminates with a legal root action distribution.
Assume the following conditions on its supported domain.
\begin{enumerate}
\item The root input is $i_q=I_q(h_{\mathrm{actual}})$. Configuration, persistent
inputs, belief construction, and initialization are fixed or factor through $i_q$.
\item All randomness and other external computational inputs have the same
joint law at equal $i_q$. In particular, fresh random bits are independent of
the actual hidden state conditional on the permitted inputs.
\item Every later operation---sampling, proposal, selection, advancement,
settlement, backup, failure handling, stopping, and final choice---uses only
these inputs and the hypotheses and records generated from them. Trusted
simulation may inspect a generated hypothesis in full.
\item The bounds and fallback make the computation finite and total. Any
external timing used to stop computation satisfies the identical-law condition
above, including its dependence on the generated computation.
\end{enumerate}
Then two actual histories with the same $I_q$ induce the same root action
distribution, namely a behavioral policy $\pi_q^{\mathrm{search}}(\cdot\mid i_q)$.
\end{proposition}
\begin{proof}
Couple two runs at equal root information using the same permitted inputs and
randomness. Their initial beliefs and records agree. Inductively, each operation
receives equal arguments, so its sampled hypothesis, successful output or
failure, and updated record agree. The same argument applies to stopping and
fallback. Finite termination gives equal final action distributions in the
coupling; averaging over its common input law proves the claim.
\end{proof}

This proposition permits substantial computation on sampled hidden state.
Its decisive restriction is that the actual hidden state contributes no extra
input beyond $i_q$. Local selector invariance is useful for the chosen search
model, but root safety follows from this broader provenance condition. Opponent
realizability still asks the cross-context question in
\cref{prop:root-conditioned}.

\begin{figure}[htbp]
\centering
\begin{tikzpicture}[alt={Root information, configuration, and randomness generate a hypothetical history inside trusted simulation. Projection sends player information and the represented menu to selection. The selected action returns to advancement. Simulation also supplies terminal, evaluator, or rollout settlement, which updates search record Z. The record feeds back into selection and supplies the final root choice. The actual hidden referee state supplies no extra input.},every node/.style={font=\small\sffamily},
 box/.style={draw=accent,fill=pale,rounded corners,text width=53mm,
 minimum height=13mm,align=center},arr/.style={-{Latex},draw=accent,thick}]
\node[box] (world) {Trusted simulation\\Sample and advance $h$};
\node[box,right=25mm of world] (project) {Player projection\\$I_r(h)$ and represented menu};
\node[box,below=17mm of world] (settle) {Settlement\\Terminal / evaluator / rollout};
\node[box,below=17mm of project] (select) {Proposal and selection\\Player input plus record $Z$};
\node[box,below=17mm of settle] (record) {Search record $Z$\\Actions, counts, and returns};
\node[box,below=17mm of select] (choice) {Final root choice\\$\rho(\cdot\mid i_q,Z)$};
\node[anchor=south west,align=left] (input) at ([yshift=12mm]world.north west)
  {Root information $i_q$, configuration $\xi$, and randomness};
\draw[arr] (world.north |- input.south)--(world.north);
\draw[arr] (world)--(project); \draw[arr] (project)--(select);
\draw[arr] (select.north west)--node[above,sloped,font=\footnotesize] {Action $a$}(world.south east);
\draw[arr] (world)--(settle); \draw[arr] (settle)--(record);
\draw[arr] (record.north east)--node[above,sloped,font=\footnotesize] {$Z$}(select.south west);
\draw[arr] (record)--(choice);
\end{tikzpicture}
\caption{Information flow in a search invocation. Selection reads projected
player information and the evolving record, then returns an action to
advancement. Full hypotheses stay within simulation and settlement. All paths
begin with permitted root inputs.}
\end{figure}

\section{Symmetry between players}
\label{sec:symmetry}

Exchanging players $0$ and $1$ carries ownership, control, private contents,
and the role of playing first with the exchange. This changes the labels while
preserving the strategic situation, including any advantage of playing first.
For the signal game, the initial family also includes the states obtained by
exchanging the sender's and guesser's player labels.

Let $\sigma$ act on players, states, actions, and observations, with
$\sigma(q)=1-q$ and $\sigma^2$ the identity on each domain. On observation
words and histories it acts componentwise, preserving order. Require
\begin{align}
\sigma(\Init)&=\Init, & \sigma(F)&=F,\\
p(\sigma(s))&=1-p(s), & A(\sigma(s))&=\sigma(A(s)),\label{eq:sym-game}\\
\nexts(\sigma(s),\sigma(a))&=\sigma(\nexts(s,a)).&&
\end{align}
The nonterminal expressions apply only where they are defined. For observations,
\begin{align}
O_{1-q}^0(\sigma(s_0))&=\sigma(O_q^0(s_0)),\\
O_{1-q}(\sigma(s),\sigma(a))&=\sigma(O_q(s,a)).
\label{eq:sym-obs}
\end{align}

\begin{proposition}[Equivariance of information and compatibility]
For every valid history $h$,
\[
I_{1-q}(\sigma(h))=\sigma(I_q(h)),\qquad
\mathcal C_{1-q}(\sigma(i_q))=\sigma(\mathcal C_q(i_q)).
\]
The corresponding game-state and epistemic-state sets transform in the same way,
where $\sigma(s,i_0,i_1)=(\sigma(s),\sigma(i_1),\sigma(i_0))$.
\end{proposition}
\begin{proof}
The game equations make $\sigma(h)$ valid. Apply the observation equations
to each term of \cref{eq:information}. For compatibility, membership on either
side follows by this identity; the reverse inclusion follows from $\sigma^2=1$.
The statements about images follow by projecting the resulting sets.
\end{proof}

\begin{example}[First player versus player zero]
An initial state may record that player $0$ plays first and owns deck $D$.
Its exchanged state records that player $1$ plays first and owns the exchanged
deck. Both states belong to the symmetric initial-state family. A convention
that always writes the first player as $0$ is a coordinate choice, not a reason
to demand that an individually labeled initial state be fixed by $\sigma$.
\end{example}

Rule symmetry can relate instances with different decks, private information,
and policies: the exchange transports each of these components to its new
owner. Requiring equal components or an exchange-invariant prior adds
restrictions on the particular instance or experiment.

\subsection{Exchanging policies}

Define the exchanged profile by
\[
\pi^\sigma_{1-q}(\sigma(a)\mid\sigma(i_q))
=\pi_q(a\mid i_q).
\]
The exchanged initial distribution is $\lambda^\sigma(\sigma(s))=\lambda(s)$.
Substituting into \cref{eq:path} gives
\[
\Prob_{\lambda^\sigma,\pi^\sigma}([\sigma(h)])
=\Prob_{\lambda,\pi}([h]).
\]
Requiring $\pi^\sigma=\pi$ is an additional assumption of symmetric behavior.
The game symmetry itself only tells us how to exchange a possibly asymmetric
profile.

\subsection{Exchanging beliefs, values, and evaluators}

Assume payoff respects player exchange: $u_{1-q}(\sigma(s))=u_q(s)$.
Transport a belief by $\sigma$ and use the exchanged continuation profile.
Pairing each continuation with its exchanged continuation preserves path
probability and the appropriately labeled payoff. Under the domains and
termination conditions of \cref{sec:value}, this gives
\begin{align*}
V_{1-q}(\sigma_{\#}b_q,\pi^\sigma)&=V_q(b_q,\pi),\\
Q_{1-q}(\sigma_{\#}b_q,\sigma(a);\pi^\sigma)&=Q_q(b_q,a;\pi).
\end{align*}
A symmetric evaluator family likewise satisfies
$e_{1-q}(\sigma(i_q))=e_q(i_q)$. This compares exchanged perspectives in
exchanged situations. Within one history, the players may have different
subjective models, as the cutoff example in \cref{sec:cutoff} illustrates.

\Needspace{28\baselineskip}
\section{Where these objects appear in MTGallium}
\label{sec:implementation}

The table points to the code representing each object and identifies what
remains to be checked. Its links follow the public source layout:
\begin{center}\code{main}\end{center}

\small
\begin{longtable}{@{}>{\raggedright\arraybackslash}p{.17\textwidth}>{\raggedright\arraybackslash}p{.25\textwidth}>{\raggedright\arraybackslash}p{.24\textwidth}>{\raggedright\arraybackslash}p{.25\textwidth}@{}}
\toprule Theory & Source representation & Mapped domain & What to check\\\midrule\endfirsthead
\toprule Theory & Source representation & Mapped domain & What to check\\\midrule\endhead
\bottomrule\endfoot
Current visible information &
\sourcepath{agent/infoset-argentum/src/main/kotlin/org/mtgallium/agent/infoset/argentum/SafeObservationProjector.kt}{SafeObservation-\newline Projector} &
Viewer-specific Argentum observation to \code{PlayerObservation}\newline\code{Snapshot} &
Does every field, reference, and ordering factor through legitimate information?\\[3pt]
Disclosures and remembered events &
\sourcepath{agent/infoset-argentum/src/main/kotlin/org/mtgallium/agent/infoset/argentum/PerspectiveHistory.kt}{PerspectiveHistory};
\sourcepath{agent/infoset-argentum/src/main/kotlin/org/mtgallium/agent/infoset/argentum/PerspectiveEventProjector.kt}{PerspectiveEvent-\newline Projector} &
Accepted choices, event projection, and visible transitions &
Do supported events preserve order, own choices, and remembered knowledge?\\[3pt]
Player information representation &
\sourcepath{agent/infoset-argentum/src/main/kotlin/org/mtgallium/agent/infoset/argentum/ArgentumSearchWorld.kt}{ArgentumSearchWorld}\newline \code{informationState} &
Snapshot, safe history, knowledge, and candidate choices &
For which tasks is this representation sufficient? Which distinctions in $I_q$ are preserved?\\[3pt]
Evaluator input and settlement origin &
\sourcepath{agent/infoset-planning/src/main/kotlin/org/mtgallium/agent/infoset/core/SearchBoundary.kt}{SearchBoundary.kt} &
Information-state and sampled-world leaf-value interfaces &
Does each call use the intended input domain, perspective, and target?\\
\end{longtable}
\normalsize

\code{PlayerObservationSnapshot} is a current snapshot, whereas the paper's observation
word accumulates disclosures through time. The source event pipeline combines
new engine events with safe references and remembered knowledge. For example,
a draw projection can include identities available to the drawer or justified
by prior revelation and tracked library order. This is a concrete way to enrich
a new event with previously derived knowledge.

The \code{argentum-board-v1} evaluator uses a sampled-world evaluation
route. Visible V2, tactical scoring, and the linear and residual evaluators
described in the \sourcepath{docs/glossary.md}{project glossary} use
information-state inputs.

\subsection{Application boundary of the deterministic model}

Applying this framework to MTGallium requires checking the engine's transitions,
disclosures, and decision points. Drawing from a fixed hidden library order
fits the deterministic model when that order is part of the state; fresh random
shuffles require a chance-transition extension. The source map identifies code
relevant to these checks, but does not establish that a complete game family
satisfies the model.
\Cref{sec:scope} lists the model ingredients that such an extension must supply.

\section{Connections to the literature}
\label{sec:related}

Extensive-form games supply histories, information sets, and behavioral
strategies. Osborne and Rubinstein~\cite{osborne1994} provide a textbook
treatment. Here the observation word is the \emph{information state}; its fiber
among acting histories is the information set. The word makes recall and
chronological updates convenient to state.

Kocsis and Szepesv\'ari~\cite[Section~2.4 and Theorem~6]{kocsis2006} analyze
UCT in specified MDP and payoff-process settings. The selection expression in
\cref{eq:ucb} illustrates the same bandit-guided allocation idea. Applying a
convergence result here would require verifying its assumptions for the chosen
hypothesis populations, adaptive continuations, and settlement rules.

Cowling, Powley, and Whitehouse~\cite[Sections~II and IV.E]{cowling2012}
develop information-set Monte Carlo search and explain strategy fusion: combining
choices across determinizations can grant a player distinctions it cannot
observe. Their Section~X addresses the searching player's private information
being implicitly available to simulated opponents, directly motivating the
cross-context question in \cref{prop:root-conditioned}. Their algorithms use
particular tree and information-set conventions; the acting-player key $K(h)$
here is an explicit choice of statistical sharing for this specification.

\section{Scope boundaries and a specification template}
\label{sec:scope}

This model leaves engine chance transitions, truly simultaneous moves,
nonterminating-play payoff conventions, intermediate rewards, discounting,
learning updates, and equilibrium computation unspecified. Each extension
needs its corresponding dynamics, information, and value conventions. Modeling
simultaneous choices through hidden commitments, for example, requires an
explicit disclosure and decision protocol.

To specify a particular search algorithm using this framework, fill in the
following items:
\begin{enumerate}
\setlength{\itemsep}{1pt}\setlength{\parsep}{0pt}
\item \textbf{Game and information domain}: initial states, legal choices,
advancement, disclosures, payoff, and supported representations.
\item \textbf{Belief model}: hypothesis object, prior or construction rule,
historical behavior assumptions, endpoint, and approximation method.
\item \textbf{Target}: payoff perspective, future policy profile or strategic
objective, and termination or cutoff meaning.
\item \textbf{Search structure}: node key, action identity, represented menu,
and successor registration under consistent hypotheses.
\item \textbf{Computation}: proposal, selection, rollout, settlement, all
budgets, failure handling, and backup conventions.
\item \textbf{Output}: root choice distribution, tie-breaking, fallback, and
whether statistics are discarded or reused, and why reuse is valid.
\end{enumerate}

The guesser's original question now has an answer under the chosen sender
model: after $\ell$, orange has probability $4/5$, so guessing orange has
expected payoff $3/5$. A short search can favor that choice while reporting a
different sample mean. We can explain the difference, track what a failed
attempt leaves out, and recognize when a search record gives the modeled
opponent an extra clue. These same links between observations, models, choices,
and returns give a larger implementation its meaning.

\appendix\clearpage

\section{Notation reference}
\label{app:notation}
\renewcommand{\arraystretch}{1.05}
\begin{longtable}{@{}p{.22\textwidth}p{.72\textwidth}@{}}
\toprule Symbol & Meaning\\\midrule\endfirsthead
\toprule Symbol & Meaning (continued)\\\midrule\endhead
\bottomrule\endfoot
$P=\{0,1\}$ & Players; labels are distinct from the role of playing first.\\
$S,\Init,F$ & Game states, allowed initial states, terminal states.\\
$p(s),A(s)$ & Acting player and nonempty finite legal action set at a nonterminal state.\\
$\nexts(s,a)$ & Deterministic advancement through automatic consequences to the next game state.\\
$h,h_j,\Hist$ & A valid finite game history, its prefix at depth $j$, and the set of all valid finite histories.\\
$\last(h)$ & The final game state of a history.\\
$\Obs_q,\Obs_q^*$ & Player $q$'s observation alphabet and finite observation words.\\
$O_q^0,O_q$ & Initial and transition observation-sequence functions.\\
$\varepsilon,\cat$ & Empty word and concatenation; history extension uses the analogous notation.\\
$I_q(h),\Info_q$ & Full player observation history and its realizable range; the theoretical information state.\\
$R_q,D_q$ & Decoders of own-choice sequence and current decision status.\\
$\Info_q^{\mathrm{act}}$ & Information states at which player $q$ must choose.\\
$A_q(i_q)$ & Legal actions decoded from an acting information state.\\
$\mathcal C_q(i_q)$ & Compatible game histories, without probabilities.\\
$\epss(h)$ & Full epistemic state $(\last(h),I_0(h),I_1(h))$.\\
$\mathcal S_q,\mathcal E_q$ & Compatible current game states and compatible full epistemic states.\\
$\sigma$ & Involutive player exchange, extended consistently to all labeled objects.\\
$\Delta(X)$ & Probability distributions on finite or countable set $X$.\\
$\pi_q,\pi$ & Behavioral policy and two-player policy profile.\\
$\lambda,[h]$ & Initial-state distribution and event that play has prefix $h$.\\
$b_q,\ f_{\#}b_q$ & Belief over compatible histories and its pushforward through $f$.\\
$u_q$ & Terminal payoff to player $q$.\\
$G_q^\pi(h)$ & Expected terminal payoff from a specified history under continuation profile $\pi$.\\
$V_q,Q_q$ & Belief value and forced-next-action value for a specified continuation profile.\\
$e_q,\widehat e_q$ & Theoretical information-state scorer and scorer on a concrete representation.\\
$r_q,\Phi_q$ & Information-state encoding and its history-based constructor $\Phi_q=r_q\circ I_q$.\\
$K(h),n$ & Acting-player information-state key and a decision node.\\
$A_{\mathrm{rep}}(n)$ & Legal actions currently represented at a search node.\\
$N,W,\widehat Q$ & Successful value-bearing visit count, backed sum, and empirical mean.\\
$(v,\kappa)$ & Settlement value and its declared origin.\\
$Z,\xi,B,\rho$ & Search record, configuration, budget, and root choice rule.\\
\end{longtable}

\section{Compact contract sheet}
\label{app:contracts}

\begingroup
\renewcommand{\labelenumi}{\textbf{C\arabic{enumi}.}}
\begin{enumerate}
\item \textbf{Decision completeness.} \emph{Model assumption; \cref{sec:game}.} Every nonterminal game state has one actor
and a nonempty finite legal action set; terminal states have no outgoing choice.
\item \textbf{Deterministic advancement.} \emph{Model assumption; \cref{sec:game}.} Legal $(s,a)$ determines exactly one
successor after finite automatic processing. No additional choice is consumed.
\item \textbf{Disclosure sufficiency.} \emph{Model assumption; \cref{sec:observation}.} $O_q(s,a)$ is determined by the chosen
state and action; any needed current visibility context is represented.
\item \textbf{Observation history.} \emph{Definition; \cref{sec:observation}.} $I_q$ is concatenated disclosures, including
initial information and correctly positioned own-choice observations.
\item \textbf{Decision awareness.} \emph{Model assumption; \cref{sec:observation}.} Equal $I_q$ values imply equal current
decision status and, when acting, equal legal action sets.
\item \textbf{Compatibility.} \emph{Consequence; \cref{sec:compatibility}.} The actual history lies in its compatibility
class; joint epistemic components share one compatible witness.
\item \textbf{Behavioral perspective.} \emph{Definition; \cref{sec:policy}.} Under a fixed policy, the action kernel
depends only on the acting player's information state.
\item \textbf{Symmetry.} \emph{Model assumption and consequence; \cref{sec:symmetry}.} Exchanging player identities commutes with advancement,
observation, exchanged policies, and exchanged payoff perspectives.
\item \textbf{Belief support.} \emph{Definition and modeling choice; \cref{sec:belief}.} Hypotheses lie within compatibility. Bayesian
claims identify the past-generation law and a finite-history endpoint.
\item \textbf{Value identity.} \emph{Definition; \cref{sec:value}.} Expected values specify belief, payoff,
continuation, and termination assumptions. A forced action is an intervention.
\item \textbf{Consistent simulation.} \emph{Implementation obligation; \cref{sec:simulation}.} The current hypothesis determines the
successor. Genuine continuation policies use only their actor's information
state. Search-conditioned selectors obey fixed-record invariance; interpreting
them as one opponent policy also requires agreement across root contexts.
\item \textbf{Settlement origins.} \emph{Algorithm convention; \cref{sec:simulation}.} Terminal payoff, heuristic scoring, and
bounded continuation retain distinct origins. Failure has no strategic value.
\item \textbf{Population-aware backup.} \emph{Algorithm convention and inference obligation; \cref{sec:estimates}.} Counts and means describe actual
value-bearing traversals. Relating them to a target accounts for the effects of
failure, selection, and changing continuations on the sampled returns.
\item \textbf{Total output or explicit partiality.} \emph{Implementation obligation; \cref{sec:algorithm}.} Root selection stays legal
and handles ties and absent values, or explicitly reports inability to choose.
\item \textbf{Representation separation.} \emph{Implementation obligation; \cref{sec:representations}.} Safety, completeness, decision
sufficiency, and execution fidelity are independently stated obligations.
\end{enumerate}
\endgroup

\begingroup
\makeatletter
\renewcommand\subsection{\Needspace{6\baselineskip}\@startsection{subsection}{2}{\z@}
  {-2.2ex \@plus -.5ex \@minus -.2ex}{1ex \@plus .2ex}
  {\normalfont\large\sffamily\bfseries\color{ink}}}
\makeatother
\section{Exercises}
\label{app:exercises}

Self-checks 1--12 revisit the main concepts; exercises 13--16 ask for a construction, proof, diagnosis, and calculation. Solutions begin in \cref{app:solutions}.

\subsection*{1. reconstructibility}
Suppose the engine draws the next card from a fixed hidden library. Does that
violate deterministic advancement? What if it shuffles before drawing?

\subsection*{2. Empty observation sequences}
A hidden opponent action changes no information available to player $q$, who
is not asked to act before or after it. Why can $O_q(s,a)=\varepsilon$? What
would change if a ``hidden action occurred'' token were appended instead?

\subsection*{3. equal board, unequal history}
Two players' current snapshots are equal in two histories, but one history
contains a remembered reveal and the other does not. Can a complete
information-state representation assign the same encoding?

\subsection*{4. a legal menu that leaks information}
Two histories have the same proposed observation word, but action $a$ is legal
in only one. What part of the model must be revisited?

\subsection*{5. Learning from one's own action}
At information state $i$, a player selects $a$ with probability $1/4$ in every
compatible history. Does learning that they selected $a$ change their belief
over those histories before any further disclosure?

\subsection*{6. forced off-policy action}
The continuation policy assigns action $a$ probability zero. Is
$Q_q(b_q,a;\pi)$ undefined?

\subsection*{7. zero-sum expectations}
Player $0$ assigns probability $0.8$ to winning; player $1$, using a different
belief, assigns probability $0.7$ to winning. With no draws, what are their
subjective values, and is zero-sum payoff violated?

\subsection*{8. means and failures}
An evaluator is exact, but the adapter fails on every hypothesis with negative
payoff. Does omitting those failures from $W$ and $N$ ensure an unbiased estimate?

\subsection*{9. symmetry and action order}
Two actions have equal root statistics. An implementation chooses whichever
raw action identifier sorts first. Is player symmetry guaranteed?

\subsection*{10. full epistemic sufficiency}
Two histories have the same $(s,I_0,I_1)$ but different unobserved past events.
Can their future payoff laws differ under the same information-state policies?

\subsection*{11. choosing the same node twice}
Does sharing an information-state key permit a simulation to pick any outgoing
successor of that graph node?

\subsection*{12. evaluator versus expected value}
A learned scorer returns $0.6$ at a nonterminal information state. List the
additional information needed before interpreting it as an expected payoff.

\subsection*{13. Construct an observation encoding}
For the signal game, construct typed initial and transition observation words for both players. Include enough information to recover own choices and current decision status. Give the guesser word after $\ell$ and explain why it is equal under both types.

\subsection*{14. Prove prefix-freeness}
Prove that an acting compatibility class is prefix-free using only own-choice recoverability and decision consistency. Explain why the same argument need not apply to a waiting class.

\subsection*{15. Diagnose two selector inputs}
Two hypotheses have equal guesser information $i_\ell$. Selector A reads their hidden type directly to choose the matching guess. Selector B reads only $(i_\ell,Z)$, but $Z$ was built from a sender-root search that knows the type. Which local and cross-context conditions can each violate?

\subsection*{16. Calculate the missing return}
Continue the trace in \cref{sec:example} with a seventh attempt that successfully guesses orange under type $\btype$. Compute both action records. Then compare the orange mean with a counterfactual execution in which attempt 6 also completed. Explain which record the actual algorithm may use.

\clearpage\section{Worked solutions}
\label{app:solutions}

\subsection*{1. reconstructibility}
A fixed library order belongs to the game state, so drawing its
next card can be deterministic even though a player does not know the identity.
A stochastic shuffle is outside the present transition model. Hiding its
random seed from the policy does not make the theoretical operation deterministic
unless one explicitly changes what the state and model include. This paper
does not make that extension.

\subsection*{2. Empty observation sequences}
Empty output leaves player history unchanged. The token would
distinguish histories according to occurrence or count of that hidden action.
It is valid only if that occurrence really is disclosed. If the transition
changes whether $q$ must act, decision awareness requires a distinguishing
observation, so the empty-output premise would no longer apply.

\subsection*{3. equal board, unequal history}
No. Completeness requires injectivity on theoretical information
states. A snapshot-only encoding may still be safe and decision-sufficient,
but it is not complete for full observation history.

\subsection*{4. a legal menu that leaks information}
Decision decoding cannot map the same word to two different
legal sets. The information state must include the legitimately disclosed menu
distinction, the action abstraction must be corrected, or the claimed
observation model is invalid. Removing $a$ from the menu only in sampled worlds
does not preserve an exact shared information-state node.

\subsection*{5. Learning from one's own action}
No: multiplying every prior mass by $1/4$ and normalizing
returns the same belief. An opponent's action may have different likelihoods
across its different private information states, so the corresponding argument
does not generally apply to observed opponent behavior.

\subsection*{6. forced off-policy action}
Not for that reason. The first action is fixed by definition;
only later choices follow $\pi$. The action must be legal throughout the root
compatibility class, and the resulting continuation must satisfy the payoff
and termination assumptions. Ordinary conditioning on a zero-probability
observed action would be a different operation.

\subsection*{7. zero-sum expectations}
Their values are $0.6$ and $0.4$. There is no violation: the
expectations use different probability distributions. Under a common
distribution, the two payoff expectations would be negatives of each other.

\subsection*{8. means and failures}
No. It avoids assigning a strategic value to failure, but
conditions the empirical population on successful execution. A claimed
all-hypothesis target still lacks the missing negative-return population.

\subsection*{9. symmetry and action order}
No. If the ordering is not preserved under player exchange,
the tie-break can change behavior under relabeling. A symmetric semantic order,
or a symmetric distribution over ties, is needed for the stated exchange law.
The same concern applies to proposal truncation.

\subsection*{10. full epistemic sufficiency}
Not under the stated sufficiency assumptions: their next action
law, successor state, and appended observations agree recursively. If the laws
do differ, a fact omitted from the tuple affects transition, disclosure, payoff,
or policy memory, so the model's sufficient-state or policy contract was false.

\subsection*{11. choosing the same node twice}
No. Edges have existential witnesses that can be different
histories. The sampled hypothesis must witness the chosen action's actual
successor. The quotient graph represents statistical sharing, not an autonomous
world simulator.

\subsection*{12. evaluator versus expected value}
At least the payoff convention, target belief construction,
continuation policies or strategic objective, supported input domain, and any
output transformation. Even then it is an estimate whose validity needs
evidence; its numerical range and input type do not establish accuracy.

\subsection*{13. Construct an observation encoding}
Use initial tokens $\operatorname{type}(\theta)$ and $\operatorname{request}(\{\ell,\mathsf r\})$ for the sender, and $\operatorname{waiting}$ for the guesser. After signal $\ell$, append $\operatorname{own}(\ell),\operatorname{waiting}$ to the sender and $\operatorname{signal}(\ell),\operatorname{request}(\{g_{\otype},g_{\btype}\})$ to the guesser. After the guess, append its own-action token to the guesser and a terminal token to both, together with declared payoff disclosure. The guesser word after $\ell$ is $[\operatorname{waiting},\operatorname{signal}(\ell),\operatorname{request}(\{g_{\otype},g_{\btype}\})]$ under both types. Filtering own-action tokens recovers choices, and the latest request, waiting, or terminal token decodes status.

\subsection*{14. Prove prefix-freeness}
If compatible $h$ is a proper prefix of compatible $h\prime$, player $q$ acts at the end of $h$ and its chosen action occurs in $h\prime$. Thus $\own_q(h\prime)$ strictly extends $\own_q(h)$. Equal information states would decode equal own-choice sequences, a contradiction. In a waiting class the intervening actions can all belong to the opponent and emit no observation to $q$, so neither the observation word nor its own-choice sequence need grow.

\subsection*{15. Diagnose two selector inputs}
A violates fixed-record invariance: changing the current hypothesis with $i_\ell,Z$ fixed changes its action. B can satisfy that local rule while violating global opponent realizability, as in \cref{prop:root-conditioned}. Both computations can still be root-perspective-safe if the hidden hypothesis and record were generated entirely from legitimate root inputs; using the actual referee type beyond those inputs would violate \cref{prop:search-safety}.

\subsection*{16. Calculate the missing return}
After the failed sixth attempt, orange has $(N,W)=(4,2)$. Attempt 7 adds $-1$, producing $(5,1,1/5)$; blue remains $(1,-1,-1)$. If attempt 6 had also completed with $-1$, orange would instead have $(6,0,0)$. The actual backup uses the former record and retains the sixth failure in attempt accounting. Counterfactual completion is an explanatory calculation, not an observed return.

\endgroup
\Needspace{10\baselineskip}
\begin{thebibliography}{9}
\addcontentsline{toc}{section}{References}
\bibitem{osborne1994}
Martin J. Osborne and Ariel Rubinstein.
\emph{A Course in Game Theory}. MIT Press, 1994.
\href{https://www.economics.utoronto.ca/osborne/cgt/INDEXR.HTM}{Authors' book page};
\href{https://books.osborne.economics.utoronto.ca/}{authorized electronic text}.

\bibitem{kocsis2006}
Levente Kocsis and Csaba Szepesv\'ari.
``Bandit based Monte-Carlo Planning.''
In \emph{Machine Learning: ECML 2006}, LNCS 4212, pp.~282--293, 2006.
\href{https://sites.ualberta.ca/~szepesva/papers/ecml06.pdf}{Author-hosted paper}.

\bibitem{cowling2012}
Peter I. Cowling, Edward J. Powley, and Daniel Whitehouse.
``Information Set Monte Carlo Tree Search.''
\emph{IEEE Transactions on Computational Intelligence and AI in Games},
4(2):120--143, 2012.
\href{https://doi.org/10.1109/TCIAIG.2012.2200894}{doi:10.1109/TCIAIG.2012.2200894}.
\href{https://eprints.whiterose.ac.uk/id/eprint/75048/1/CowlingPowleyWhitehouse2012.pdf}{University repository copy}.
\end{thebibliography}

\end{document}
